% Options for packages loaded elsewhere
\PassOptionsToPackage{unicode}{hyperref}
\PassOptionsToPackage{hyphens}{url}
\PassOptionsToPackage{dvipsnames,svgnames,x11names}{xcolor}
%
\documentclass[
  10pt,
]{article}

\usepackage{amsmath,amssymb}
\usepackage{iftex}
\ifPDFTeX
  \usepackage[T1]{fontenc}
  \usepackage[utf8]{inputenc}
  \usepackage{textcomp} % provide euro and other symbols
\else % if luatex or xetex
  \usepackage{unicode-math}
  \defaultfontfeatures{Scale=MatchLowercase}
  \defaultfontfeatures[\rmfamily]{Ligatures=TeX,Scale=1}
\fi
\usepackage{lmodern}
\ifPDFTeX\else  
    % xetex/luatex font selection
    \setmainfont[]{Arial}
\fi
% Use upquote if available, for straight quotes in verbatim environments
\IfFileExists{upquote.sty}{\usepackage{upquote}}{}
\IfFileExists{microtype.sty}{% use microtype if available
  \usepackage[]{microtype}
  \UseMicrotypeSet[protrusion]{basicmath} % disable protrusion for tt fonts
}{}
\makeatletter
\@ifundefined{KOMAClassName}{% if non-KOMA class
  \IfFileExists{parskip.sty}{%
    \usepackage{parskip}
  }{% else
    \setlength{\parindent}{0pt}
    \setlength{\parskip}{6pt plus 2pt minus 1pt}}
}{% if KOMA class
  \KOMAoptions{parskip=half}}
\makeatother
\usepackage{xcolor}
\usepackage[a4paper]{geometry}
\setlength{\emergencystretch}{3em} % prevent overfull lines
\setcounter{secnumdepth}{5}
% Make \paragraph and \subparagraph free-standing
\makeatletter
\ifx\paragraph\undefined\else
  \let\oldparagraph\paragraph
  \renewcommand{\paragraph}{
    \@ifstar
      \xxxParagraphStar
      \xxxParagraphNoStar
  }
  \newcommand{\xxxParagraphStar}[1]{\oldparagraph*{#1}\mbox{}}
  \newcommand{\xxxParagraphNoStar}[1]{\oldparagraph{#1}\mbox{}}
\fi
\ifx\subparagraph\undefined\else
  \let\oldsubparagraph\subparagraph
  \renewcommand{\subparagraph}{
    \@ifstar
      \xxxSubParagraphStar
      \xxxSubParagraphNoStar
  }
  \newcommand{\xxxSubParagraphStar}[1]{\oldsubparagraph*{#1}\mbox{}}
  \newcommand{\xxxSubParagraphNoStar}[1]{\oldsubparagraph{#1}\mbox{}}
\fi
\makeatother

\usepackage{color}
\usepackage{fancyvrb}
\newcommand{\VerbBar}{|}
\newcommand{\VERB}{\Verb[commandchars=\\\{\}]}
\DefineVerbatimEnvironment{Highlighting}{Verbatim}{commandchars=\\\{\}}
% Add ',fontsize=\small' for more characters per line
\usepackage{framed}
\definecolor{shadecolor}{RGB}{248,248,248}
\newenvironment{Shaded}{\begin{snugshade}}{\end{snugshade}}
\newcommand{\AlertTok}[1]{\textcolor[rgb]{0.94,0.16,0.16}{#1}}
\newcommand{\AnnotationTok}[1]{\textcolor[rgb]{0.56,0.35,0.01}{\textbf{\textit{#1}}}}
\newcommand{\AttributeTok}[1]{\textcolor[rgb]{0.13,0.29,0.53}{#1}}
\newcommand{\BaseNTok}[1]{\textcolor[rgb]{0.00,0.00,0.81}{#1}}
\newcommand{\BuiltInTok}[1]{#1}
\newcommand{\CharTok}[1]{\textcolor[rgb]{0.31,0.60,0.02}{#1}}
\newcommand{\CommentTok}[1]{\textcolor[rgb]{0.56,0.35,0.01}{\textit{#1}}}
\newcommand{\CommentVarTok}[1]{\textcolor[rgb]{0.56,0.35,0.01}{\textbf{\textit{#1}}}}
\newcommand{\ConstantTok}[1]{\textcolor[rgb]{0.56,0.35,0.01}{#1}}
\newcommand{\ControlFlowTok}[1]{\textcolor[rgb]{0.13,0.29,0.53}{\textbf{#1}}}
\newcommand{\DataTypeTok}[1]{\textcolor[rgb]{0.13,0.29,0.53}{#1}}
\newcommand{\DecValTok}[1]{\textcolor[rgb]{0.00,0.00,0.81}{#1}}
\newcommand{\DocumentationTok}[1]{\textcolor[rgb]{0.56,0.35,0.01}{\textbf{\textit{#1}}}}
\newcommand{\ErrorTok}[1]{\textcolor[rgb]{0.64,0.00,0.00}{\textbf{#1}}}
\newcommand{\ExtensionTok}[1]{#1}
\newcommand{\FloatTok}[1]{\textcolor[rgb]{0.00,0.00,0.81}{#1}}
\newcommand{\FunctionTok}[1]{\textcolor[rgb]{0.13,0.29,0.53}{\textbf{#1}}}
\newcommand{\ImportTok}[1]{#1}
\newcommand{\InformationTok}[1]{\textcolor[rgb]{0.56,0.35,0.01}{\textbf{\textit{#1}}}}
\newcommand{\KeywordTok}[1]{\textcolor[rgb]{0.13,0.29,0.53}{\textbf{#1}}}
\newcommand{\NormalTok}[1]{#1}
\newcommand{\OperatorTok}[1]{\textcolor[rgb]{0.81,0.36,0.00}{\textbf{#1}}}
\newcommand{\OtherTok}[1]{\textcolor[rgb]{0.56,0.35,0.01}{#1}}
\newcommand{\PreprocessorTok}[1]{\textcolor[rgb]{0.56,0.35,0.01}{\textit{#1}}}
\newcommand{\RegionMarkerTok}[1]{#1}
\newcommand{\SpecialCharTok}[1]{\textcolor[rgb]{0.81,0.36,0.00}{\textbf{#1}}}
\newcommand{\SpecialStringTok}[1]{\textcolor[rgb]{0.31,0.60,0.02}{#1}}
\newcommand{\StringTok}[1]{\textcolor[rgb]{0.31,0.60,0.02}{#1}}
\newcommand{\VariableTok}[1]{\textcolor[rgb]{0.00,0.00,0.00}{#1}}
\newcommand{\VerbatimStringTok}[1]{\textcolor[rgb]{0.31,0.60,0.02}{#1}}
\newcommand{\WarningTok}[1]{\textcolor[rgb]{0.56,0.35,0.01}{\textbf{\textit{#1}}}}

\providecommand{\tightlist}{%
  \setlength{\itemsep}{0pt}\setlength{\parskip}{0pt}}\usepackage{longtable,booktabs,array}
\usepackage{calc} % for calculating minipage widths
% Correct order of tables after \paragraph or \subparagraph
\usepackage{etoolbox}
\makeatletter
\patchcmd\longtable{\par}{\if@noskipsec\mbox{}\fi\par}{}{}
\makeatother
% Allow footnotes in longtable head/foot
\IfFileExists{footnotehyper.sty}{\usepackage{footnotehyper}}{\usepackage{footnote}}
\makesavenoteenv{longtable}
\usepackage{graphicx}
\makeatletter
\newsavebox\pandoc@box
\newcommand*\pandocbounded[1]{% scales image to fit in text height/width
  \sbox\pandoc@box{#1}%
  \Gscale@div\@tempa{\textheight}{\dimexpr\ht\pandoc@box+\dp\pandoc@box\relax}%
  \Gscale@div\@tempb{\linewidth}{\wd\pandoc@box}%
  \ifdim\@tempb\p@<\@tempa\p@\let\@tempa\@tempb\fi% select the smaller of both
  \ifdim\@tempa\p@<\p@\scalebox{\@tempa}{\usebox\pandoc@box}%
  \else\usebox{\pandoc@box}%
  \fi%
}
% Set default figure placement to htbp
\def\fps@figure{htbp}
\makeatother

\makeatletter
\@ifpackageloaded{caption}{}{\usepackage{caption}}
\AtBeginDocument{%
\ifdefined\contentsname
  \renewcommand*\contentsname{Table of contents}
\else
  \newcommand\contentsname{Table of contents}
\fi
\ifdefined\listfigurename
  \renewcommand*\listfigurename{List of Figures}
\else
  \newcommand\listfigurename{List of Figures}
\fi
\ifdefined\listtablename
  \renewcommand*\listtablename{List of Tables}
\else
  \newcommand\listtablename{List of Tables}
\fi
\ifdefined\figurename
  \renewcommand*\figurename{Figure}
\else
  \newcommand\figurename{Figure}
\fi
\ifdefined\tablename
  \renewcommand*\tablename{Table}
\else
  \newcommand\tablename{Table}
\fi
}
\@ifpackageloaded{float}{}{\usepackage{float}}
\floatstyle{ruled}
\@ifundefined{c@chapter}{\newfloat{codelisting}{h}{lop}}{\newfloat{codelisting}{h}{lop}[chapter]}
\floatname{codelisting}{Listing}
\newcommand*\listoflistings{\listof{codelisting}{List of Listings}}
\makeatother
\makeatletter
\makeatother
\makeatletter
\@ifpackageloaded{caption}{}{\usepackage{caption}}
\@ifpackageloaded{subcaption}{}{\usepackage{subcaption}}
\makeatother

\usepackage{bookmark}

\IfFileExists{xurl.sty}{\usepackage{xurl}}{} % add URL line breaks if available
\urlstyle{same} % disable monospaced font for URLs
\hypersetup{
  pdftitle={04\_Mutant\_feedbacks},
  pdfauthor={Henry Janse van Rensburg},
  colorlinks=true,
  linkcolor={blue},
  filecolor={Maroon},
  citecolor={Blue},
  urlcolor={Blue},
  pdfcreator={LaTeX via pandoc}}


\title{04\_Mutant\_feedbacks}
\author{Henry Janse van Rensburg}
\date{2025-01-31}

\begin{document}
\maketitle

\renewcommand*\contentsname{Table of contents}
{
\hypersetup{linkcolor=}
\setcounter{tocdepth}{2}
\tableofcontents
}

\section{Background}\label{background}

We conducted this experiment to study the effect of \emph{mmf1} mutants,
and mutants affected in ETI signalling pathways on BX-growth feedbacks.
We repeated screening of \emph{mmf1} mutants three times

\section{Materials and methods}\label{materials-and-methods}

We grew Arabidopsis plants on BXplus and BXminus soils and phenotyped
them for rosette area and weighed them at the end of the experiment. We
extracted the phenotypes using ARADEEPOPSIS (@hüther2020).

\textbf{Experiments:}

\begin{itemize}
\tightlist
\item
  \textbf{\emph{mmf1}} \textbf{mutants}:
\end{itemize}

1- Screening of \emph{mmf1-1} on BS04 soil batch

2- Screening of \emph{mmf1-1} and \emph{mmf1-2} BS04 soil batch

3- Screening of \emph{mmf1-1} and \emph{mmf1-2} BS09 soil batch

\begin{itemize}
\tightlist
\item
  \textbf{ETI mutants}:
\end{itemize}

4- Screening of \emph{pad4-1}, \emph{eds1-12} and \emph{sag101-3} on
BS07 soil batch

5- Screening of \emph{pad4-1}, \emph{eds1-12} and \emph{sag101-3} on
BS07 soil batch

\begin{itemize}
\tightlist
\item
  \textbf{Purpose}: We performed this study to confirm the involvement
  of MMF1 in BX-feedbacks.
\end{itemize}

\section{Important results}\label{important-results}

\emph{mmf1} mutants had a significant effect on the growth of
Arabidopsis plants on BXplus and BXminus soils. They showed no growth
feedback. ETI mutants on the other hand showed a significant growth
feedback on BXplus soils.

\begin{Shaded}
\begin{Highlighting}[]
\DocumentationTok{\#\#\# **1. Setup and Load Libraries**}

\CommentTok{\# Load necessary libraries}
\FunctionTok{library}\NormalTok{(ggplot2)}
\FunctionTok{library}\NormalTok{(ggbeeswarm)}
\FunctionTok{library}\NormalTok{(tidyverse)}
\FunctionTok{library}\NormalTok{(ggtext)}
\FunctionTok{library}\NormalTok{(tidyr)}
\FunctionTok{library}\NormalTok{(ggpubr)}
\FunctionTok{library}\NormalTok{(pander)}
\FunctionTok{library}\NormalTok{(cowplot)}
\FunctionTok{library}\NormalTok{(dplyr)}
\FunctionTok{library}\NormalTok{(knitr)}
\FunctionTok{library}\NormalTok{(writexl)}
\FunctionTok{library}\NormalTok{(openxlsx)}
\FunctionTok{library}\NormalTok{(readxl)}
\end{Highlighting}
\end{Shaded}

\section{Read Phenotype Tables}\label{read-phenotype-tables}

\begin{Shaded}
\begin{Highlighting}[]
\CommentTok{\# Read phenotype tables}
\CommentTok{\# Exp 1}
\NormalTok{phenotypes\_E1 }\OtherTok{\textless{}{-}} \FunctionTok{read\_excel}\NormalTok{(}\StringTok{"Input/Dataset4.xlsx"}\NormalTok{, }\AttributeTok{sheet =} \StringTok{"ExpIII\_mmf1\_phenotypes"}\NormalTok{)}
\NormalTok{weight\_E1 }\OtherTok{\textless{}{-}} \FunctionTok{read\_excel}\NormalTok{(}\StringTok{"Input/Dataset4.xlsx"}\NormalTok{, }\AttributeTok{sheet =} \StringTok{"ExpIII\_mmf1\_weight"}\NormalTok{)}
 
\CommentTok{\#Exp2}
\NormalTok{phenotypes\_E2 }\OtherTok{\textless{}{-}} \FunctionTok{read\_excel}\NormalTok{(}\StringTok{"Input/Dataset4.xlsx"}\NormalTok{, }\AttributeTok{sheet =} \StringTok{"ExpIV\_mmf1\_phenotypes"}\NormalTok{)}
\NormalTok{weight\_E2 }\OtherTok{\textless{}{-}} \FunctionTok{read\_excel}\NormalTok{(}\StringTok{"Input/Dataset4.xlsx"}\NormalTok{, }\AttributeTok{sheet =} \StringTok{"ExpIV\_mmf1\_weight"}\NormalTok{)}
  
\CommentTok{\#Exp3}
\NormalTok{phenotypes\_E3 }\OtherTok{\textless{}{-}} \FunctionTok{read\_excel}\NormalTok{(}\StringTok{"Input/Dataset4.xlsx"}\NormalTok{, }\AttributeTok{sheet =} \StringTok{"ExpV\_mmf1\_phenotypes"}\NormalTok{)}
  
\NormalTok{weight\_E3 }\OtherTok{\textless{}{-}} \FunctionTok{read\_excel}\NormalTok{(}\StringTok{"Input/Dataset4.xlsx"}\NormalTok{, }\AttributeTok{sheet =} \StringTok{"ExpV\_mmf1\_weight"}\NormalTok{)}
  
\CommentTok{\#removed\_plants\_E3 \textless{}{-} read.csv("Input/E3\_removed\_plants.csv")}
\CommentTok{\# Purpose: Load the phenotype data and mutant screen fresh weight data into R.}
\end{Highlighting}
\end{Shaded}

\section{Function to Add Replicate Number to
Graph}\label{function-to-add-replicate-number-to-graph}

\begin{Shaded}
\begin{Highlighting}[]
\CommentTok{\# Function to display the number of replicates below each bar}
\NormalTok{give.n }\OtherTok{\textless{}{-}} \ControlFlowTok{function}\NormalTok{(x)\{}
  \FunctionTok{return}\NormalTok{(}\FunctionTok{c}\NormalTok{(}\AttributeTok{y =} \FunctionTok{median}\NormalTok{(x) }\SpecialCharTok{*} \SpecialCharTok{{-}}\FloatTok{0.005}\NormalTok{, }\AttributeTok{label =} \FunctionTok{length}\NormalTok{(x)))}
\NormalTok{\}}

\CommentTok{\# Function to make colored facets at the bottom}
\NormalTok{element\_textbox\_highlight }\OtherTok{\textless{}{-}} \ControlFlowTok{function}\NormalTok{(..., }
                                      \AttributeTok{hi.labels =} \ConstantTok{NULL}\NormalTok{, }\AttributeTok{hi.fill =} \ConstantTok{NULL}\NormalTok{,}
                                      \AttributeTok{hi.col =} \ConstantTok{NULL}\NormalTok{, }\AttributeTok{hi.box.col =} \ConstantTok{NULL}\NormalTok{,}
                                      \AttributeTok{hi.labels2 =} \ConstantTok{NULL}\NormalTok{, }\AttributeTok{hi.fill2 =} \ConstantTok{NULL}\NormalTok{,}
                                      \AttributeTok{hi.col2 =} \ConstantTok{NULL}\NormalTok{, }\AttributeTok{hi.box.col2 =} \ConstantTok{NULL}\NormalTok{) \{}
  \FunctionTok{structure}\NormalTok{(}
    \FunctionTok{c}\NormalTok{(}\FunctionTok{element\_textbox}\NormalTok{(...),}
      \FunctionTok{list}\NormalTok{(}\AttributeTok{hi.labels =}\NormalTok{ hi.labels, }\AttributeTok{hi.fill =}\NormalTok{ hi.fill, }\AttributeTok{hi.col =}\NormalTok{ hi.col, }
           \AttributeTok{hi.box.col =}\NormalTok{ hi.box.col,}
           \AttributeTok{hi.labels2 =}\NormalTok{ hi.labels2, }\AttributeTok{hi.fill2 =}\NormalTok{ hi.fill2, }\AttributeTok{hi.col2 =}\NormalTok{ hi.col2, }
           \AttributeTok{hi.box.col2 =}\NormalTok{ hi.box.col2)}
\NormalTok{    ),}
    \AttributeTok{class =} \FunctionTok{c}\NormalTok{(}\StringTok{"element\_textbox\_highlight"}\NormalTok{, }\StringTok{"element\_textbox"}\NormalTok{, }\StringTok{"element\_text"}\NormalTok{, }\StringTok{"element"}\NormalTok{,}
              \StringTok{"element\_textbox\_highlight"}\NormalTok{, }\StringTok{"element\_textbox"}\NormalTok{, }\StringTok{"element\_text"}\NormalTok{, }\StringTok{"element"}\NormalTok{)}
\NormalTok{  )}
\NormalTok{\}}

\NormalTok{element\_grob.element\_textbox\_highlight }\OtherTok{\textless{}{-}} \ControlFlowTok{function}\NormalTok{(element, }\AttributeTok{label =} \StringTok{""}\NormalTok{, ...) \{}
  \ControlFlowTok{if}\NormalTok{ (label }\SpecialCharTok{\%in\%}\NormalTok{ element}\SpecialCharTok{$}\NormalTok{hi.labels) \{}
\NormalTok{    element}\SpecialCharTok{$}\NormalTok{fill }\OtherTok{\textless{}{-}}\NormalTok{ element}\SpecialCharTok{$}\NormalTok{hi.fill }\SpecialCharTok{\%||\%}\NormalTok{ element}\SpecialCharTok{$}\NormalTok{fill}
\NormalTok{    element}\SpecialCharTok{$}\NormalTok{colour }\OtherTok{\textless{}{-}}\NormalTok{ element}\SpecialCharTok{$}\NormalTok{hi.col }\SpecialCharTok{\%||\%}\NormalTok{ element}\SpecialCharTok{$}\NormalTok{colour}
\NormalTok{    element}\SpecialCharTok{$}\NormalTok{box.colour }\OtherTok{\textless{}{-}}\NormalTok{ element}\SpecialCharTok{$}\NormalTok{hi.box.col }\SpecialCharTok{\%||\%}\NormalTok{ element}\SpecialCharTok{$}\NormalTok{box.colour}
\NormalTok{  \}}
  \ControlFlowTok{if}\NormalTok{ (label }\SpecialCharTok{\%in\%}\NormalTok{ element}\SpecialCharTok{$}\NormalTok{hi.labels2) \{}
\NormalTok{    element}\SpecialCharTok{$}\NormalTok{fill }\OtherTok{\textless{}{-}}\NormalTok{ element}\SpecialCharTok{$}\NormalTok{hi.fill2 }\SpecialCharTok{\%||\%}\NormalTok{ element}\SpecialCharTok{$}\NormalTok{fill}
\NormalTok{    element}\SpecialCharTok{$}\NormalTok{colour }\OtherTok{\textless{}{-}}\NormalTok{ element}\SpecialCharTok{$}\NormalTok{hi.col2 }\SpecialCharTok{\%||\%}\NormalTok{ element}\SpecialCharTok{$}\NormalTok{colour}
\NormalTok{    element}\SpecialCharTok{$}\NormalTok{box.colour }\OtherTok{\textless{}{-}}\NormalTok{ element}\SpecialCharTok{$}\NormalTok{hi.box.col2 }\SpecialCharTok{\%||\%}\NormalTok{ element}\SpecialCharTok{$}\NormalTok{box.colour}
\NormalTok{  \}}
  \FunctionTok{NextMethod}\NormalTok{()}
\NormalTok{\}}
\CommentTok{\# Purpose: Define functions for visualizing the number of replicates in a graph and customizing text box elements.}
\end{Highlighting}
\end{Shaded}

\section{\texorpdfstring{Experiment III: Screening of \emph{mmf1-1} on
BS04 soil
batch}{Experiment III: Screening of mmf1-1 on BS04 soil batch}}\label{experiment-iii-screening-of-mmf1-1-on-bs04-soil-batch}

\subsection{Modify and Verify Data}\label{modify-and-verify-data}

\begin{Shaded}
\begin{Highlighting}[]
\CommentTok{\# Change the levels of the \textquotesingle{}line\textquotesingle{} column}
\NormalTok{weight\_E1 }\OtherTok{\textless{}{-}}\NormalTok{ weight\_E1 }\SpecialCharTok{\%\textgreater{}\%}
  \FunctionTok{mutate}\NormalTok{(}\AttributeTok{line =} \FunctionTok{factor}\NormalTok{(line, }\AttributeTok{levels =} \FunctionTok{c}\NormalTok{(}\StringTok{"WT"}\NormalTok{, }\StringTok{"mmf1{-}1"}\NormalTok{)))}

\CommentTok{\#number of replicates}
\FunctionTok{kable}\NormalTok{(}
  \FunctionTok{table}\NormalTok{(weight\_E1}\SpecialCharTok{$}\NormalTok{BX, weight\_E1}\SpecialCharTok{$}\NormalTok{line), }
  \AttributeTok{caption =} \StringTok{"Number of replicates (E{-}III)"}
\NormalTok{)}
\end{Highlighting}
\end{Shaded}

\begin{longtable}[]{@{}lrr@{}}
\caption{Number of replicates (E-III)}\tabularnewline
\toprule\noalign{}
& WT & mmf1-1 \\
\midrule\noalign{}
\endfirsthead
\toprule\noalign{}
& WT & mmf1-1 \\
\midrule\noalign{}
\endhead
\bottomrule\noalign{}
\endlastfoot
minus & 14 & 14 \\
plus & 15 & 15 \\
\end{longtable}

\begin{Shaded}
\begin{Highlighting}[]
\CommentTok{\# Purpose: Relevel the \textquotesingle{}line\textquotesingle{} factor and filter out rows with missing fresh weight data.}
\end{Highlighting}
\end{Shaded}

\subsection{Plot Fresh Weight}\label{plot-fresh-weight}

\begin{Shaded}
\begin{Highlighting}[]
\CommentTok{\# Generate the boxplot}
\NormalTok{fresh\_weight\_E1 }\OtherTok{\textless{}{-}}\NormalTok{ weight\_E1 }\SpecialCharTok{\%\textgreater{}\%}
  \FunctionTok{ggplot}\NormalTok{(}\FunctionTok{aes}\NormalTok{(}\AttributeTok{x =}\NormalTok{ line, }\AttributeTok{y =}\NormalTok{ fresh\_weight, }\AttributeTok{fill =}\NormalTok{ BX)) }\SpecialCharTok{+} 
  \FunctionTok{geom\_boxplot}\NormalTok{(}\AttributeTok{position =} \FunctionTok{position\_dodge}\NormalTok{(}\AttributeTok{width =} \FloatTok{0.85}\NormalTok{), }\AttributeTok{alpha =} \FloatTok{0.6}\NormalTok{, }
               \AttributeTok{color =} \StringTok{"grey35"}\NormalTok{,}
               \AttributeTok{outlier.colour =} \ConstantTok{NA}\NormalTok{, }\AttributeTok{size =} \FloatTok{0.5}\NormalTok{) }\SpecialCharTok{+} 
\NormalTok{  ggbeeswarm}\SpecialCharTok{::}\FunctionTok{geom\_quasirandom}\NormalTok{(}\AttributeTok{dodge.width =} \FloatTok{0.85}\NormalTok{, }\AttributeTok{size =} \DecValTok{1}\NormalTok{, }\AttributeTok{shape =} \DecValTok{21}\NormalTok{, }
                               \AttributeTok{color =} \StringTok{"black"}\NormalTok{,}
                               \AttributeTok{width =} \FloatTok{0.1}\NormalTok{, }\AttributeTok{alpha =} \FloatTok{0.4}\NormalTok{, }\AttributeTok{show.legend =} \ConstantTok{FALSE}\NormalTok{) }\SpecialCharTok{+}
  \FunctionTok{stat\_summary}\NormalTok{(}\AttributeTok{fun.data =}\NormalTok{ mean\_se, }\AttributeTok{geom =} \StringTok{"errorbar"}\NormalTok{, }\AttributeTok{linewidth =} \FloatTok{0.5}\NormalTok{, }
               \AttributeTok{width =} \FloatTok{0.4}\NormalTok{, }\AttributeTok{position =} \FunctionTok{position\_dodge}\NormalTok{(}\AttributeTok{width =} \FloatTok{0.85}\NormalTok{)) }\SpecialCharTok{+} 
  \FunctionTok{stat\_compare\_means}\NormalTok{(}\AttributeTok{method =} \StringTok{"t.test"}\NormalTok{, }\AttributeTok{paired =} \ConstantTok{FALSE}\NormalTok{, }\AttributeTok{size =} \DecValTok{2}\NormalTok{, }
                     \AttributeTok{label.y =} \FloatTok{0.09}\NormalTok{, }\AttributeTok{label.x =} \DecValTok{1}\NormalTok{) }\SpecialCharTok{+}
  \FunctionTok{theme\_bw}\NormalTok{() }\SpecialCharTok{+}
  \FunctionTok{scale\_fill\_manual}\NormalTok{(}\AttributeTok{name =} \StringTok{"Conditioning"}\NormalTok{, }\AttributeTok{values =} \FunctionTok{c}\NormalTok{(}\StringTok{"seagreen"}\NormalTok{, }\StringTok{"gold2"}\NormalTok{)) }\SpecialCharTok{+} 
  \FunctionTok{coord\_cartesian}\NormalTok{(}\AttributeTok{ylim =} \FunctionTok{c}\NormalTok{(}\DecValTok{0}\NormalTok{, }\FloatTok{0.10}\NormalTok{)) }\SpecialCharTok{+}
  \FunctionTok{ggtitle}\NormalTok{(}\StringTok{"E1 {-} Fresh Weight"}\NormalTok{) }\SpecialCharTok{+}  \CommentTok{\# Add a title}
  \FunctionTok{ylab}\NormalTok{(}\StringTok{"Fresh Weight"}\NormalTok{) }\SpecialCharTok{+} 
  \FunctionTok{scale\_x\_discrete}\NormalTok{(}\AttributeTok{labels =} \FunctionTok{c}\NormalTok{(}\StringTok{"WT"}\NormalTok{, }\FunctionTok{expression}\NormalTok{(}\FunctionTok{italic}\NormalTok{(}\StringTok{"mmf1{-}2"}\NormalTok{)))) }\SpecialCharTok{+} 
  \FunctionTok{theme}\NormalTok{(}
    \AttributeTok{plot.title =} \FunctionTok{element\_text}\NormalTok{(}\AttributeTok{size =} \DecValTok{7}\NormalTok{, }\AttributeTok{hjust =} \DecValTok{0}\NormalTok{),  }\CommentTok{\# Move title to the left}
    \AttributeTok{axis.title.x =} \FunctionTok{element\_text}\NormalTok{(}\AttributeTok{size =} \DecValTok{7}\NormalTok{),}
    \AttributeTok{axis.title.y =} \FunctionTok{element\_text}\NormalTok{(}\AttributeTok{size =} \DecValTok{7}\NormalTok{),}
    \AttributeTok{axis.text.x =} \FunctionTok{element\_text}\NormalTok{(}\AttributeTok{size =} \DecValTok{7}\NormalTok{),}
    \AttributeTok{axis.text.y =} \FunctionTok{element\_text}\NormalTok{(}\AttributeTok{size =} \DecValTok{7}\NormalTok{),}
    \AttributeTok{legend.title =} \FunctionTok{element\_text}\NormalTok{(}\AttributeTok{size =} \DecValTok{7}\NormalTok{),}
    \AttributeTok{legend.text =} \FunctionTok{element\_text}\NormalTok{(}\AttributeTok{size =} \DecValTok{7}\NormalTok{),}
    \AttributeTok{strip.text =} \FunctionTok{element\_text}\NormalTok{(}\AttributeTok{size =} \DecValTok{7}\NormalTok{),  }\CommentTok{\# Facet strip text size}
    \AttributeTok{legend.position =} \StringTok{"bottom"}
\NormalTok{  )}

\CommentTok{\# Display the plot}
\FunctionTok{print}\NormalTok{(fresh\_weight\_E1)}
\CommentTok{\# Purpose: Generate and display a boxplot for fresh weight data in MMF1 mutants.}
\end{Highlighting}
\end{Shaded}

\begin{figure}[H]

\centering{

\pandocbounded{\includegraphics[keepaspectratio]{04_Mutant_feedbacks_files/figure-pdf/fig-mmf1-mutants-Fw-EIII-1.pdf}}

}

\caption{\label{fig-mmf1-mutants-Fw-EIII}Feedbacks of mmf1 mutants in
manuscript experiments (E-III).}

\end{figure}%

\subsection{Plot Plant Area}\label{plot-plant-area}

\begin{Shaded}
\begin{Highlighting}[]
\NormalTok{phenotypes\_E1 }\OtherTok{\textless{}{-}}\NormalTok{ phenotypes\_E1 }\SpecialCharTok{\%\textgreater{}\%}
  \FunctionTok{mutate}\NormalTok{(}\AttributeTok{line =} \FunctionTok{factor}\NormalTok{(line, }\AttributeTok{levels =} \FunctionTok{c}\NormalTok{(}\StringTok{"WT"}\NormalTok{, }\StringTok{"mmf1{-}1"}\NormalTok{)))}


\CommentTok{\#number of replicates}
\FunctionTok{kable}\NormalTok{(}
  \FunctionTok{table}\NormalTok{(phenotypes\_E1}\SpecialCharTok{$}\NormalTok{BX, phenotypes\_E1}\SpecialCharTok{$}\NormalTok{line), }
  \AttributeTok{caption =} \StringTok{"Number of replicates (E{-}III)"}
\NormalTok{)}
\end{Highlighting}
\end{Shaded}

\begin{longtable}[]{@{}lrr@{}}
\caption{Number of replicates (E-III)}\tabularnewline
\toprule\noalign{}
& WT & mmf1-1 \\
\midrule\noalign{}
\endfirsthead
\toprule\noalign{}
& WT & mmf1-1 \\
\midrule\noalign{}
\endhead
\bottomrule\noalign{}
\endlastfoot
minus & 14 & 14 \\
plus & 15 & 15 \\
\end{longtable}

\begin{Shaded}
\begin{Highlighting}[]
\CommentTok{\# Generate the boxplot}
\NormalTok{plant\_area\_E1 }\OtherTok{\textless{}{-}}\NormalTok{ phenotypes\_E1 }\SpecialCharTok{\%\textgreater{}\%}
  \FunctionTok{ggplot}\NormalTok{(}\FunctionTok{aes}\NormalTok{(}\AttributeTok{x =}\NormalTok{ line, }\AttributeTok{y =}\NormalTok{ plant\_region\_area, }\AttributeTok{fill =}\NormalTok{ BX)) }\SpecialCharTok{+} 
  \FunctionTok{geom\_boxplot}\NormalTok{(}\AttributeTok{position =} \FunctionTok{position\_dodge}\NormalTok{(}\AttributeTok{width =} \FloatTok{0.85}\NormalTok{), }\AttributeTok{alpha =} \FloatTok{0.6}\NormalTok{, }
               \AttributeTok{color =} \StringTok{"grey35"}\NormalTok{, }\AttributeTok{outlier.colour =} \ConstantTok{NA}\NormalTok{, }\AttributeTok{size =} \FloatTok{0.5}\NormalTok{) }\SpecialCharTok{+} 
\NormalTok{  ggbeeswarm}\SpecialCharTok{::}\FunctionTok{geom\_quasirandom}\NormalTok{(}\AttributeTok{dodge.width =} \FloatTok{0.85}\NormalTok{, }\AttributeTok{size =} \DecValTok{1}\NormalTok{, }\AttributeTok{shape =} \DecValTok{21}\NormalTok{, }
                               \AttributeTok{color =} \StringTok{"black"}\NormalTok{, }\AttributeTok{width =} \FloatTok{0.1}\NormalTok{, }\AttributeTok{alpha =} \FloatTok{0.4}\NormalTok{, }
                               \AttributeTok{show.legend =} \ConstantTok{FALSE}\NormalTok{) }\SpecialCharTok{+}
  \FunctionTok{stat\_summary}\NormalTok{(}\AttributeTok{fun.data =}\NormalTok{ mean\_se, }\AttributeTok{geom =} \StringTok{"errorbar"}\NormalTok{, }\AttributeTok{linewidth =} \FloatTok{0.5}\NormalTok{, }
               \AttributeTok{width =} \FloatTok{0.4}\NormalTok{, }\AttributeTok{position =} \FunctionTok{position\_dodge}\NormalTok{(}\AttributeTok{width =} \FloatTok{0.85}\NormalTok{)) }\SpecialCharTok{+} 
  \FunctionTok{stat\_compare\_means}\NormalTok{(}\AttributeTok{method =} \StringTok{"t.test"}\NormalTok{, }\AttributeTok{paired =} \ConstantTok{FALSE}\NormalTok{, }\AttributeTok{size =} \DecValTok{2}\NormalTok{, }
                     \AttributeTok{label.y =} \DecValTok{900000}\NormalTok{, }\AttributeTok{label.x =} \DecValTok{1}\NormalTok{) }\SpecialCharTok{+}
  \FunctionTok{theme\_bw}\NormalTok{() }\SpecialCharTok{+}
  \FunctionTok{scale\_fill\_manual}\NormalTok{(}\AttributeTok{name =} \StringTok{"Conditioning"}\NormalTok{, }\AttributeTok{values =} \FunctionTok{c}\NormalTok{(}\StringTok{"seagreen"}\NormalTok{, }\StringTok{"gold2"}\NormalTok{)) }\SpecialCharTok{+} 
  \FunctionTok{coord\_cartesian}\NormalTok{(}\AttributeTok{ylim =} \FunctionTok{c}\NormalTok{(}\DecValTok{0}\NormalTok{, }\DecValTok{1000000}\NormalTok{)) }\SpecialCharTok{+}
  \FunctionTok{ggtitle}\NormalTok{(}\StringTok{"E1 {-} Plant Area"}\NormalTok{) }\SpecialCharTok{+}  \CommentTok{\# Add a title}
  \FunctionTok{ylab}\NormalTok{(}\StringTok{"Plant Area"}\NormalTok{) }\SpecialCharTok{+} 
  \FunctionTok{scale\_x\_discrete}\NormalTok{(}\AttributeTok{labels =} \FunctionTok{c}\NormalTok{(}\StringTok{"WT"}\NormalTok{, }\FunctionTok{expression}\NormalTok{(}\FunctionTok{italic}\NormalTok{(}\StringTok{"mmf1{-}1"}\NormalTok{)))) }\SpecialCharTok{+}  \CommentTok{\# Adjust x{-}axis labels}
  \FunctionTok{theme}\NormalTok{(}
    \AttributeTok{plot.title =} \FunctionTok{element\_text}\NormalTok{(}\AttributeTok{size =} \DecValTok{7}\NormalTok{, }\AttributeTok{hjust =} \DecValTok{0}\NormalTok{),  }\CommentTok{\# Title aligned to the left}
    \AttributeTok{axis.title.x =} \FunctionTok{element\_blank}\NormalTok{(),  }\CommentTok{\# Optional: Remove x{-}axis title}
    \AttributeTok{axis.title.y =} \FunctionTok{element\_text}\NormalTok{(}\AttributeTok{size =} \DecValTok{7}\NormalTok{),}
    \AttributeTok{axis.text.x =} \FunctionTok{element\_text}\NormalTok{(}\AttributeTok{size =} \DecValTok{7}\NormalTok{),}
    \AttributeTok{axis.text.y =} \FunctionTok{element\_text}\NormalTok{(}\AttributeTok{size =} \DecValTok{7}\NormalTok{),}
    \AttributeTok{legend.title =} \FunctionTok{element\_text}\NormalTok{(}\AttributeTok{size =} \DecValTok{7}\NormalTok{),}
    \AttributeTok{legend.text =} \FunctionTok{element\_text}\NormalTok{(}\AttributeTok{size =} \DecValTok{7}\NormalTok{),}
    \AttributeTok{strip.text =} \FunctionTok{element\_text}\NormalTok{(}\AttributeTok{size =} \DecValTok{7}\NormalTok{),  }\CommentTok{\# Facet strip text size}
    \AttributeTok{legend.position =} \StringTok{"bottom"}
\NormalTok{  )}

\CommentTok{\# Display the plot}
\FunctionTok{print}\NormalTok{(plant\_area\_E1)}
\CommentTok{\# Purpose: Generate and display a boxplot for plant area data in MMF1 mutants.}
\end{Highlighting}
\end{Shaded}

\begin{figure}[H]

\centering{

\pandocbounded{\includegraphics[keepaspectratio]{04_Mutant_feedbacks_files/figure-pdf/fig-mmf1-mutants-area-EIII-1.pdf}}

}

\caption{\label{fig-mmf1-mutants-area-EIII}Feedbacks of mmf1 mutants in
repeat experiments (E-III).}

\end{figure}%

\subsection{Statistics}\label{statistics}

\begin{Shaded}
\begin{Highlighting}[]
\CommentTok{\# Perform ANOVA for MMF1 mutants}
\CommentTok{\# Fit the ANOVA model for fresh weight}
\CommentTok{\# Fit the ANOVA model}
\NormalTok{MMF1\_anova\_model\_Fw\_E1 }\OtherTok{\textless{}{-}} \FunctionTok{aov}\NormalTok{(fresh\_weight }\SpecialCharTok{\textasciitilde{}}\NormalTok{ line }\SpecialCharTok{*}\NormalTok{ BX, }\AttributeTok{data =}\NormalTok{ weight\_E1)}

\CommentTok{\# Extract the ANOVA table and format it with kable}
\NormalTok{anova\_table }\OtherTok{\textless{}{-}} \FunctionTok{summary}\NormalTok{(MMF1\_anova\_model\_Fw\_E1)[[}\DecValTok{1}\NormalTok{]]  }\CommentTok{\# Extract the ANOVA table as a data frame}

\CommentTok{\# Display the ANOVA table using kable}
\FunctionTok{library}\NormalTok{(knitr)}
\FunctionTok{kable}\NormalTok{(anova\_table, }\AttributeTok{caption =} \StringTok{"ANOVA Table for Fresh Weight (E{-}III)"}\NormalTok{)}
\end{Highlighting}
\end{Shaded}

\begin{longtable}[]{@{}lrrrrr@{}}
\caption{ANOVA Table for Fresh Weight (E-III)}\tabularnewline
\toprule\noalign{}
& Df & Sum Sq & Mean Sq & F value & Pr(\textgreater F) \\
\midrule\noalign{}
\endfirsthead
\toprule\noalign{}
& Df & Sum Sq & Mean Sq & F value & Pr(\textgreater F) \\
\midrule\noalign{}
\endhead
\bottomrule\noalign{}
\endlastfoot
line & 1 & 0.0001128 & 0.0001128 & 0.5779622 & 0.4504188 \\
BX & 1 & 0.0020235 & 0.0020235 & 10.3665983 & 0.0021736 \\
line:BX & 1 & 0.0002511 & 0.0002511 & 1.2862520 & 0.2617495 \\
Residuals & 54 & 0.0105404 & 0.0001952 & NA & NA \\
\end{longtable}

\begin{Shaded}
\begin{Highlighting}[]
\CommentTok{\# Fit the ANOVA model for plant region area}
\NormalTok{MMF1\_anova\_model\_rosettearea\_E1 }\OtherTok{\textless{}{-}} \FunctionTok{aov}\NormalTok{(plant\_region\_area }\SpecialCharTok{\textasciitilde{}}\NormalTok{ line }\SpecialCharTok{*}\NormalTok{ BX, }
                                       \AttributeTok{data =}\NormalTok{ phenotypes\_E1)}

\CommentTok{\# Extract the ANOVA table and format it with kable}
\NormalTok{anova\_table }\OtherTok{\textless{}{-}} \FunctionTok{summary}\NormalTok{(MMF1\_anova\_model\_rosettearea\_E1)[[}\DecValTok{1}\NormalTok{]]  }\CommentTok{\# Extract the ANOVA table as a data frame}

\CommentTok{\# Display the ANOVA table using kable}
\FunctionTok{kable}\NormalTok{(anova\_table, }\AttributeTok{caption =} \StringTok{"ANOVA Table for Plant Region Area (E{-}III)"}\NormalTok{)}
\end{Highlighting}
\end{Shaded}

\begin{longtable}[]{@{}lrrrrr@{}}
\caption{ANOVA Table for Plant Region Area (E-III)}\tabularnewline
\toprule\noalign{}
& Df & Sum Sq & Mean Sq & F value & Pr(\textgreater F) \\
\midrule\noalign{}
\endfirsthead
\toprule\noalign{}
& Df & Sum Sq & Mean Sq & F value & Pr(\textgreater F) \\
\midrule\noalign{}
\endhead
\bottomrule\noalign{}
\endlastfoot
line & 1 & 1.283010e+10 & 12830104212 & 0.6702395 & 0.4165665 \\
BX & 1 & 2.276057e+11 & 227605673691 & 11.8900284 & 0.0011014 \\
line:BX & 1 & 4.708198e+10 & 47081980430 & 2.4595436 & 0.1226545 \\
Residuals & 54 & 1.033699e+12 & 19142567756 & NA & NA \\
\end{longtable}

\begin{Shaded}
\begin{Highlighting}[]
\CommentTok{\# Purpose: Perform ANOVA to analyze the effects of line and BX on fresh weight and plant region area in MMF1 mutants.}
\end{Highlighting}
\end{Shaded}

\begin{itemize}
\tightlist
\item
  \textbf{Purpose}: We performed ANOVA to analyze the effects of line
  and BX on fresh weight and plant region area in MMF1 mutants.
\end{itemize}

\subsubsection{Perform ANOVA and FDR Correction for fresh
weight}\label{perform-anova-and-fdr-correction-for-fresh-weight}

\begin{Shaded}
\begin{Highlighting}[]
\CommentTok{\# Perform ANOVA and FDR correction for fresh weight}
\NormalTok{stat.test\_FW }\OtherTok{\textless{}{-}}\NormalTok{ weight\_E1 }\SpecialCharTok{\%\textgreater{}\%}
  \FunctionTok{group\_by}\NormalTok{(line) }\SpecialCharTok{\%\textgreater{}\%}
  \FunctionTok{do}\NormalTok{(\{}
\NormalTok{    model }\OtherTok{\textless{}{-}} \FunctionTok{aov}\NormalTok{(fresh\_weight }\SpecialCharTok{\textasciitilde{}}\NormalTok{ BX, }\AttributeTok{data =}\NormalTok{ .)}
\NormalTok{    anova\_results }\OtherTok{\textless{}{-}} \FunctionTok{summary}\NormalTok{(model)[[}\DecValTok{1}\NormalTok{]]}
\NormalTok{    p\_value }\OtherTok{\textless{}{-}}\NormalTok{ anova\_results[[}\StringTok{"Pr(\textgreater{}F)"}\NormalTok{]][}\DecValTok{1}\NormalTok{]}
    \FunctionTok{data.frame}\NormalTok{(}\AttributeTok{p\_value =}\NormalTok{ p\_value)}
\NormalTok{  \}) }\SpecialCharTok{\%\textgreater{}\%}
  \FunctionTok{ungroup}\NormalTok{() }\SpecialCharTok{\%\textgreater{}\%}
  \FunctionTok{mutate}\NormalTok{(}
    \AttributeTok{p.adj =} \FunctionTok{p.adjust}\NormalTok{(p\_value, }\AttributeTok{method =} \StringTok{"fdr"}\NormalTok{),}
    \AttributeTok{significance =} \FunctionTok{cut}\NormalTok{(p.adj, }
                       \AttributeTok{breaks =} \FunctionTok{c}\NormalTok{(}\DecValTok{0}\NormalTok{, }\FloatTok{1e{-}04}\NormalTok{, }\FloatTok{0.001}\NormalTok{, }\FloatTok{0.01}\NormalTok{, }\FloatTok{0.05}\NormalTok{, }\DecValTok{1}\NormalTok{),}
                       \AttributeTok{labels =} \FunctionTok{c}\NormalTok{(}\StringTok{"****"}\NormalTok{, }\StringTok{"***"}\NormalTok{, }\StringTok{"**"}\NormalTok{, }\StringTok{"*"}\NormalTok{, }\StringTok{"ns"}\NormalTok{))}
\NormalTok{  )}

\CommentTok{\# Display the result using kable}
\FunctionTok{kable}\NormalTok{(stat.test\_FW, }\AttributeTok{caption =} \StringTok{"Statistical Test Results for Fresh Weight (E{-}III)"}\NormalTok{)}
\end{Highlighting}
\end{Shaded}

\begin{longtable}[]{@{}lrrl@{}}
\caption{Statistical Test Results for Fresh Weight
(E-III)}\tabularnewline
\toprule\noalign{}
line & p\_value & p.adj & significance \\
\midrule\noalign{}
\endfirsthead
\toprule\noalign{}
line & p\_value & p.adj & significance \\
\midrule\noalign{}
\endhead
\bottomrule\noalign{}
\endlastfoot
WT & 0.0159009 & 0.0318019 & * \\
mmf1-1 & 0.0608325 & 0.0608325 & ns \\
\end{longtable}

\subsubsection{\texorpdfstring{\textbf{Perform ANOVA and FDR correction
for Rosette
area}}{Perform ANOVA and FDR correction for Rosette area}}\label{perform-anova-and-fdr-correction-for-rosette-area}

\begin{Shaded}
\begin{Highlighting}[]
\NormalTok{stat.test\_region\_area }\OtherTok{\textless{}{-}}\NormalTok{ phenotypes\_E1 }\SpecialCharTok{\%\textgreater{}\%}
  \FunctionTok{group\_by}\NormalTok{(line) }\SpecialCharTok{\%\textgreater{}\%}
  \FunctionTok{do}\NormalTok{(\{}
\NormalTok{    model }\OtherTok{\textless{}{-}} \FunctionTok{aov}\NormalTok{(plant\_region\_area }\SpecialCharTok{\textasciitilde{}}\NormalTok{ BX, }\AttributeTok{data =}\NormalTok{ .)}
\NormalTok{    anova\_results }\OtherTok{\textless{}{-}} \FunctionTok{summary}\NormalTok{(model)[[}\DecValTok{1}\NormalTok{]]}
\NormalTok{    p\_value }\OtherTok{\textless{}{-}}\NormalTok{ anova\_results[[}\StringTok{"Pr(\textgreater{}F)"}\NormalTok{]][}\DecValTok{1}\NormalTok{]}
    \FunctionTok{data.frame}\NormalTok{(}\AttributeTok{p\_value =}\NormalTok{ p\_value)}
\NormalTok{  \}) }\SpecialCharTok{\%\textgreater{}\%}
  \FunctionTok{ungroup}\NormalTok{() }\SpecialCharTok{\%\textgreater{}\%}
  \FunctionTok{mutate}\NormalTok{(}
    \AttributeTok{p.adj =} \FunctionTok{p.adjust}\NormalTok{(p\_value, }\AttributeTok{method =} \StringTok{"fdr"}\NormalTok{),}
    \AttributeTok{significance =} \FunctionTok{cut}\NormalTok{(}
\NormalTok{      p.adj,}
      \AttributeTok{breaks =} \FunctionTok{c}\NormalTok{(}\DecValTok{0}\NormalTok{, }\FloatTok{1e{-}04}\NormalTok{, }\FloatTok{0.001}\NormalTok{, }\FloatTok{0.01}\NormalTok{, }\FloatTok{0.05}\NormalTok{, }\DecValTok{1}\NormalTok{),}
      \AttributeTok{labels =} \FunctionTok{c}\NormalTok{(}\StringTok{"****"}\NormalTok{, }\StringTok{"***"}\NormalTok{, }\StringTok{"**"}\NormalTok{, }\StringTok{"*"}\NormalTok{, }\StringTok{"ns"}\NormalTok{)}
\NormalTok{    )}
\NormalTok{  )}

\CommentTok{\# Display the results using kable}
\FunctionTok{library}\NormalTok{(knitr)}
\FunctionTok{kable}\NormalTok{(stat.test\_region\_area, }\AttributeTok{caption =} \StringTok{"Statistical Test Results for Plant Region Area (E{-}III)"}\NormalTok{)}
\end{Highlighting}
\end{Shaded}

\begin{longtable}[]{@{}lrrl@{}}
\caption{Statistical Test Results for Plant Region Area
(E-III)}\tabularnewline
\toprule\noalign{}
line & p\_value & p.adj & significance \\
\midrule\noalign{}
\endfirsthead
\toprule\noalign{}
line & p\_value & p.adj & significance \\
\midrule\noalign{}
\endhead
\bottomrule\noalign{}
\endlastfoot
WT & 0.0039416 & 0.0078831 & ** \\
mmf1-1 & 0.1323721 & 0.1323721 & ns \\
\end{longtable}

\section{\texorpdfstring{Experiment IV: Screening of \emph{mmf1-1} and
\emph{mmf1-2} on BS04 soil
batch}{Experiment IV: Screening of mmf1-1 and mmf1-2 on BS04 soil batch}}\label{experiment-iv-screening-of-mmf1-1-and-mmf1-2-on-bs04-soil-batch}

Experiment that was used for microbiome analysis and RNAseq.

\subsection{Modify and Verify Data}\label{modify-and-verify-data-1}

\begin{Shaded}
\begin{Highlighting}[]
\CommentTok{\# Change the levels of the \textquotesingle{}line\textquotesingle{} column}
\NormalTok{weight\_E2 }\OtherTok{\textless{}{-}}\NormalTok{ weight\_E2 }\SpecialCharTok{\%\textgreater{}\%}
  \FunctionTok{mutate}\NormalTok{(}\AttributeTok{line =} \FunctionTok{factor}\NormalTok{(line, }\AttributeTok{levels =} \FunctionTok{c}\NormalTok{(}\StringTok{"WT"}\NormalTok{, }\StringTok{"mmf1{-}1"}\NormalTok{, }\StringTok{"mmf1{-}2"}\NormalTok{)))}

\CommentTok{\# Filter out rows with NA in the fresh\_weight column}
\NormalTok{weight\_E2 }\OtherTok{\textless{}{-}}\NormalTok{ weight\_E2 }\SpecialCharTok{\%\textgreater{}\%}
  \FunctionTok{filter}\NormalTok{(}\SpecialCharTok{!}\FunctionTok{is.na}\NormalTok{(fresh\_weight))}

\CommentTok{\#number of replicates}
\FunctionTok{kable}\NormalTok{(}
  \FunctionTok{table}\NormalTok{(weight\_E2}\SpecialCharTok{$}\NormalTok{BX, weight\_E2}\SpecialCharTok{$}\NormalTok{line), }
  \AttributeTok{caption =} \StringTok{"Number of replicates (E{-}IV)"}
\NormalTok{)}
\end{Highlighting}
\end{Shaded}

\begin{longtable}[]{@{}lrrr@{}}
\caption{Number of replicates (E-IV)}\tabularnewline
\toprule\noalign{}
& WT & mmf1-1 & mmf1-2 \\
\midrule\noalign{}
\endfirsthead
\toprule\noalign{}
& WT & mmf1-1 & mmf1-2 \\
\midrule\noalign{}
\endhead
\bottomrule\noalign{}
\endlastfoot
minus & 15 & 15 & 15 \\
plus & 15 & 14 & 14 \\
\end{longtable}

\begin{itemize}
\tightlist
\item
  \textbf{Purpose}: We re-leveled the `line' factor and filtered out
  rows with missing fresh weight data.
\end{itemize}

\subsection{Plot Fresh weight}\label{plot-fresh-weight-1}

\begin{Shaded}
\begin{Highlighting}[]
\CommentTok{\# Generate the boxplot}
\NormalTok{fresh\_weight\_E2 }\OtherTok{\textless{}{-}}\NormalTok{ weight\_E2 }\SpecialCharTok{\%\textgreater{}\%}
  \FunctionTok{ggplot}\NormalTok{(}\FunctionTok{aes}\NormalTok{(}\AttributeTok{x =}\NormalTok{ line, }\AttributeTok{y =}\NormalTok{ fresh\_weight, }\AttributeTok{fill =}\NormalTok{ BX)) }\SpecialCharTok{+} 
  \FunctionTok{geom\_boxplot}\NormalTok{(}\AttributeTok{position =} \FunctionTok{position\_dodge}\NormalTok{(}\AttributeTok{width =} \FloatTok{0.85}\NormalTok{), }\AttributeTok{alpha =} \FloatTok{0.6}\NormalTok{, }
               \AttributeTok{color =} \StringTok{"grey35"}\NormalTok{, }\AttributeTok{outlier.colour =} \ConstantTok{NA}\NormalTok{, }\AttributeTok{size =} \FloatTok{0.5}\NormalTok{) }\SpecialCharTok{+} 
\NormalTok{  ggbeeswarm}\SpecialCharTok{::}\FunctionTok{geom\_quasirandom}\NormalTok{(}\AttributeTok{dodge.width =} \FloatTok{0.85}\NormalTok{, }\AttributeTok{size =} \DecValTok{1}\NormalTok{, }\AttributeTok{shape =} \DecValTok{21}\NormalTok{, }
                               \AttributeTok{color =} \StringTok{"black"}\NormalTok{, }\AttributeTok{width =} \FloatTok{0.1}\NormalTok{, }\AttributeTok{alpha =} \FloatTok{0.4}\NormalTok{, }
                               \AttributeTok{show.legend =} \ConstantTok{FALSE}\NormalTok{) }\SpecialCharTok{+}
  \FunctionTok{stat\_summary}\NormalTok{(}\AttributeTok{fun.data =}\NormalTok{ mean\_se, }\AttributeTok{geom =} \StringTok{"errorbar"}\NormalTok{, }\AttributeTok{linewidth =} \FloatTok{0.5}\NormalTok{, }
               \AttributeTok{width =} \FloatTok{0.4}\NormalTok{, }\AttributeTok{position =} \FunctionTok{position\_dodge}\NormalTok{(}\AttributeTok{width =} \FloatTok{0.85}\NormalTok{)) }\SpecialCharTok{+} 
  \FunctionTok{stat\_compare\_means}\NormalTok{(}\AttributeTok{method =} \StringTok{"t.test"}\NormalTok{, }\AttributeTok{paired =} \ConstantTok{FALSE}\NormalTok{, }\AttributeTok{size =} \DecValTok{2}\NormalTok{, }
                     \AttributeTok{label.y =} \FloatTok{0.350}\NormalTok{, }\AttributeTok{label.x =} \DecValTok{1}\NormalTok{) }\SpecialCharTok{+}
  \FunctionTok{theme\_bw}\NormalTok{() }\SpecialCharTok{+}
  \FunctionTok{scale\_fill\_manual}\NormalTok{(}\AttributeTok{name =} \StringTok{"Conditioning"}\NormalTok{, }\AttributeTok{values =} \FunctionTok{c}\NormalTok{(}\StringTok{"seagreen"}\NormalTok{, }\StringTok{"gold2"}\NormalTok{)) }\SpecialCharTok{+} 
  \FunctionTok{coord\_cartesian}\NormalTok{(}\AttributeTok{ylim =} \FunctionTok{c}\NormalTok{(}\DecValTok{0}\NormalTok{, }\FloatTok{0.40}\NormalTok{)) }\SpecialCharTok{+}
  \FunctionTok{ggtitle}\NormalTok{(}\StringTok{"E2 {-} Fresh Weight Analysis"}\NormalTok{) }\SpecialCharTok{+}  \CommentTok{\# Add a title}
  \FunctionTok{ylab}\NormalTok{(}\StringTok{"Fresh Weight"}\NormalTok{) }\SpecialCharTok{+} 
  \FunctionTok{scale\_x\_discrete}\NormalTok{(}\AttributeTok{labels =} \FunctionTok{c}\NormalTok{(}\StringTok{"WT"}\NormalTok{, }\FunctionTok{expression}\NormalTok{(}\FunctionTok{italic}\NormalTok{(}\StringTok{"mmf1{-}1"}\NormalTok{)), }
                              \FunctionTok{expression}\NormalTok{(}\FunctionTok{italic}\NormalTok{(}\StringTok{"mmf1{-}2"}\NormalTok{)))) }\SpecialCharTok{+}  \CommentTok{\# Adjust x{-}axis labels}
  \FunctionTok{theme}\NormalTok{(}
    \AttributeTok{plot.title =} \FunctionTok{element\_text}\NormalTok{(}\AttributeTok{size =} \DecValTok{7}\NormalTok{, }\AttributeTok{hjust =} \DecValTok{0}\NormalTok{),  }\CommentTok{\# Title aligned to the left}
    \AttributeTok{axis.title.x =} \FunctionTok{element\_blank}\NormalTok{(),  }\CommentTok{\# Optional: Remove x{-}axis title}
    \AttributeTok{axis.title.y =} \FunctionTok{element\_text}\NormalTok{(}\AttributeTok{size =} \DecValTok{7}\NormalTok{),}
    \AttributeTok{axis.text.x =} \FunctionTok{element\_text}\NormalTok{(}\AttributeTok{size =} \DecValTok{7}\NormalTok{),}
    \AttributeTok{axis.text.y =} \FunctionTok{element\_text}\NormalTok{(}\AttributeTok{size =} \DecValTok{7}\NormalTok{),}
    \AttributeTok{legend.title =} \FunctionTok{element\_text}\NormalTok{(}\AttributeTok{size =} \DecValTok{7}\NormalTok{),}
    \AttributeTok{legend.text =} \FunctionTok{element\_text}\NormalTok{(}\AttributeTok{size =} \DecValTok{7}\NormalTok{),}
    \AttributeTok{strip.text =} \FunctionTok{element\_text}\NormalTok{(}\AttributeTok{size =} \DecValTok{7}\NormalTok{),  }\CommentTok{\# Facet strip text size}
    \AttributeTok{legend.position =} \StringTok{"bottom"}
\NormalTok{  )}

\CommentTok{\# Display the plot}
\FunctionTok{print}\NormalTok{(fresh\_weight\_E2)}
\CommentTok{\# Purpose: Generate and display a boxplot for fresh weight data in MMF1 mutants.}
\end{Highlighting}
\end{Shaded}

\begin{figure}[H]

\centering{

\pandocbounded{\includegraphics[keepaspectratio]{04_Mutant_feedbacks_files/figure-pdf/fig-mmf1-mutants-Fw-EIV-1.pdf}}

}

\caption{\label{fig-mmf1-mutants-Fw-EIV}Feedbacks of mmf1 mutants in
manuscript experiments (E-IV).}

\end{figure}%

\begin{itemize}
\tightlist
\item
  \textbf{Purpose}: We generated a boxplot to visualize the fresh weight
  data in MMF1 mutants.
\end{itemize}

\subsection{Plot Plant Area}\label{plot-plant-area-1}

\begin{longtable}[]{@{}lrrr@{}}
\caption{Number of replicates (E-IV)}\tabularnewline
\toprule\noalign{}
& WT & mmf1-1 & mmf1-2 \\
\midrule\noalign{}
\endfirsthead
\toprule\noalign{}
& WT & mmf1-1 & mmf1-2 \\
\midrule\noalign{}
\endhead
\bottomrule\noalign{}
\endlastfoot
minus & 15 & 15 & 15 \\
plus & 15 & 14 & 14 \\
\end{longtable}

\subsection{Plot Plant Area for MMF1
Mutants}\label{plot-plant-area-for-mmf1-mutants}

\begin{Shaded}
\begin{Highlighting}[]
\CommentTok{\# Generate the boxplot}
\NormalTok{plant\_area\_E2 }\OtherTok{\textless{}{-}}\NormalTok{ phenotypes\_E2 }\SpecialCharTok{\%\textgreater{}\%}
  \FunctionTok{ggplot}\NormalTok{(}\FunctionTok{aes}\NormalTok{(}\AttributeTok{x =}\NormalTok{ line, }\AttributeTok{y =}\NormalTok{ plant\_region\_area, }\AttributeTok{fill =}\NormalTok{ BX)) }\SpecialCharTok{+} 
  \FunctionTok{geom\_boxplot}\NormalTok{(}\AttributeTok{position =} \FunctionTok{position\_dodge}\NormalTok{(}\AttributeTok{width =} \FloatTok{0.85}\NormalTok{), }\AttributeTok{alpha =} \FloatTok{0.6}\NormalTok{, }
               \AttributeTok{color =} \StringTok{"grey35"}\NormalTok{, }\AttributeTok{outlier.colour =} \ConstantTok{NA}\NormalTok{, }\AttributeTok{size =} \FloatTok{0.5}\NormalTok{) }\SpecialCharTok{+} 
\NormalTok{  ggbeeswarm}\SpecialCharTok{::}\FunctionTok{geom\_quasirandom}\NormalTok{(}\AttributeTok{dodge.width =} \FloatTok{0.85}\NormalTok{, }\AttributeTok{size =} \DecValTok{1}\NormalTok{, }\AttributeTok{shape =} \DecValTok{21}\NormalTok{, }
                               \AttributeTok{color =} \StringTok{"black"}\NormalTok{,}
                               \AttributeTok{width =} \FloatTok{0.1}\NormalTok{, }\AttributeTok{alpha =} \FloatTok{0.4}\NormalTok{, }\AttributeTok{show.legend =} \ConstantTok{FALSE}\NormalTok{) }\SpecialCharTok{+}
  \FunctionTok{stat\_summary}\NormalTok{(}\AttributeTok{fun.data =}\NormalTok{ mean\_se, }\AttributeTok{geom =} \StringTok{"errorbar"}\NormalTok{, }\AttributeTok{linewidth =} \FloatTok{0.5}\NormalTok{, }\AttributeTok{width =} \FloatTok{0.4}\NormalTok{,}
               \AttributeTok{position =} \FunctionTok{position\_dodge}\NormalTok{(}\AttributeTok{width =} \FloatTok{0.85}\NormalTok{)) }\SpecialCharTok{+} 
  \FunctionTok{stat\_compare\_means}\NormalTok{(}\AttributeTok{method =} \StringTok{"t.test"}\NormalTok{, }\AttributeTok{paired =} \ConstantTok{FALSE}\NormalTok{, }\AttributeTok{size =} \DecValTok{2}\NormalTok{, }
                     \AttributeTok{label.y =} \DecValTok{800000}\NormalTok{, }\AttributeTok{label.x =} \DecValTok{1}\NormalTok{) }\SpecialCharTok{+}
  \FunctionTok{theme\_bw}\NormalTok{() }\SpecialCharTok{+}
  \FunctionTok{scale\_fill\_manual}\NormalTok{(}\AttributeTok{name =} \StringTok{"Conditioning"}\NormalTok{, }\AttributeTok{values =} \FunctionTok{c}\NormalTok{(}\StringTok{"seagreen"}\NormalTok{, }\StringTok{"gold2"}\NormalTok{)) }\SpecialCharTok{+} 
  \FunctionTok{coord\_cartesian}\NormalTok{(}\AttributeTok{ylim =} \FunctionTok{c}\NormalTok{(}\DecValTok{0}\NormalTok{, }\DecValTok{900000}\NormalTok{)) }\SpecialCharTok{+}
  \FunctionTok{ggtitle}\NormalTok{(}\StringTok{"E2 {-} Plant Area"}\NormalTok{) }\SpecialCharTok{+}  \CommentTok{\# Add a title}
  \FunctionTok{ylab}\NormalTok{(}\StringTok{"Plant Area"}\NormalTok{) }\SpecialCharTok{+} 
  \FunctionTok{scale\_x\_discrete}\NormalTok{(}\AttributeTok{labels =} \FunctionTok{c}\NormalTok{(}\StringTok{"WT"}\NormalTok{, }\FunctionTok{expression}\NormalTok{(}\FunctionTok{italic}\NormalTok{(}\StringTok{"mmf1{-}1"}\NormalTok{)), }
                              \FunctionTok{expression}\NormalTok{(}\FunctionTok{italic}\NormalTok{(}\StringTok{"mmf1{-}2"}\NormalTok{)))) }\SpecialCharTok{+}  \CommentTok{\# Adjust x{-}axis labels}
  \FunctionTok{theme}\NormalTok{(}
    \AttributeTok{plot.title =} \FunctionTok{element\_text}\NormalTok{(}\AttributeTok{size =} \DecValTok{7}\NormalTok{, }\AttributeTok{hjust =} \DecValTok{0}\NormalTok{),  }\CommentTok{\# Title aligned to the left}
    \AttributeTok{axis.title.x =} \FunctionTok{element\_blank}\NormalTok{(),  }\CommentTok{\# Optional: Remove x{-}axis title}
    \AttributeTok{axis.title.y =} \FunctionTok{element\_text}\NormalTok{(}\AttributeTok{size =} \DecValTok{7}\NormalTok{),}
    \AttributeTok{axis.text.x =} \FunctionTok{element\_text}\NormalTok{(}\AttributeTok{size =} \DecValTok{7}\NormalTok{),}
    \AttributeTok{axis.text.y =} \FunctionTok{element\_text}\NormalTok{(}\AttributeTok{size =} \DecValTok{7}\NormalTok{),}
    \AttributeTok{legend.title =} \FunctionTok{element\_text}\NormalTok{(}\AttributeTok{size =} \DecValTok{7}\NormalTok{),}
    \AttributeTok{legend.text =} \FunctionTok{element\_text}\NormalTok{(}\AttributeTok{size =} \DecValTok{7}\NormalTok{),}
    \AttributeTok{strip.text =} \FunctionTok{element\_text}\NormalTok{(}\AttributeTok{size =} \DecValTok{7}\NormalTok{),  }\CommentTok{\# Facet strip text size}
    \AttributeTok{legend.position =} \StringTok{"bottom"}
\NormalTok{  )}

\CommentTok{\# Display the plot}
\FunctionTok{print}\NormalTok{(plant\_area\_E2)}
\CommentTok{\# Purpose: Generate and display a boxplot for plant area data in MMF1 mutants.}
\end{Highlighting}
\end{Shaded}

\begin{figure}[H]

\centering{

\pandocbounded{\includegraphics[keepaspectratio]{04_Mutant_feedbacks_files/figure-pdf/fig-mmf1-mutants-area-EIV-1.pdf}}

}

\caption{\label{fig-mmf1-mutants-area-EIV}Feedbacks of mmf1 mutants in
manuscript experiments (E-IV).}

\end{figure}%

\begin{itemize}
\tightlist
\item
  \textbf{Purpose}: We generated a boxplot to visualize the plant area
  data in MMF1 mutants.
\end{itemize}

\subsection{Statistics}\label{statistics-1}

\begin{longtable}[]{@{}lrrrrr@{}}
\caption{ANOVA Table for Fresh Weight (E-IV)}\tabularnewline
\toprule\noalign{}
& Df & Sum Sq & Mean Sq & F value & Pr(\textgreater F) \\
\midrule\noalign{}
\endfirsthead
\toprule\noalign{}
& Df & Sum Sq & Mean Sq & F value & Pr(\textgreater F) \\
\midrule\noalign{}
\endhead
\bottomrule\noalign{}
\endlastfoot
line & 2 & 0.0764945 & 0.0382473 & 89.646602 & 0.0000000 \\
BX & 1 & 0.0034435 & 0.0034435 & 8.071188 & 0.0056715 \\
line:BX & 2 & 0.0014299 & 0.0007149 & 1.675694 & 0.1935235 \\
Residuals & 82 & 0.0349849 & 0.0004266 & NA & NA \\
\end{longtable}

\begin{longtable}[]{@{}lrrrrr@{}}
\caption{ANOVA Table for Plant Region Area (E-IV)}\tabularnewline
\toprule\noalign{}
& Df & Sum Sq & Mean Sq & F value & Pr(\textgreater F) \\
\midrule\noalign{}
\endfirsthead
\toprule\noalign{}
& Df & Sum Sq & Mean Sq & F value & Pr(\textgreater F) \\
\midrule\noalign{}
\endhead
\bottomrule\noalign{}
\endlastfoot
line & 2 & 304039816844 & 152019908422 & 21.457669 & 0.0000000 \\
BX & 1 & 29340301711 & 29340301711 & 4.141395 & 0.0450772 \\
line:BX & 2 & 30629767064 & 15314883532 & 2.161702 & 0.1216458 \\
Residuals & 82 & 580940661226 & 7084642210 & NA & NA \\
\end{longtable}

\begin{itemize}
\tightlist
\item
  \textbf{Purpose}: We performed ANOVA to analyze the effects of line
  and BX on fresh weight and plant region area in MMF1 mutants.
\end{itemize}

\subsubsection{Perform ANOVA and FDR Correction for fresh
weight}\label{perform-anova-and-fdr-correction-for-fresh-weight-1}

\begin{Shaded}
\begin{Highlighting}[]
\CommentTok{\# Perform ANOVA and FDR correction for fresh weight}
\NormalTok{stat.test\_FW }\OtherTok{\textless{}{-}}\NormalTok{ weight\_E2 }\SpecialCharTok{\%\textgreater{}\%}
  \FunctionTok{group\_by}\NormalTok{(line) }\SpecialCharTok{\%\textgreater{}\%}
  \FunctionTok{do}\NormalTok{(\{}
\NormalTok{    model }\OtherTok{\textless{}{-}} \FunctionTok{aov}\NormalTok{(fresh\_weight }\SpecialCharTok{\textasciitilde{}}\NormalTok{ BX, }\AttributeTok{data =}\NormalTok{ .)}
\NormalTok{    anova\_results }\OtherTok{\textless{}{-}} \FunctionTok{summary}\NormalTok{(model)[[}\DecValTok{1}\NormalTok{]]}
\NormalTok{    p\_value }\OtherTok{\textless{}{-}}\NormalTok{ anova\_results[[}\StringTok{"Pr(\textgreater{}F)"}\NormalTok{]][}\DecValTok{1}\NormalTok{]}
    \FunctionTok{data.frame}\NormalTok{(}\AttributeTok{p\_value =}\NormalTok{ p\_value)}
\NormalTok{  \}) }\SpecialCharTok{\%\textgreater{}\%}
  \FunctionTok{ungroup}\NormalTok{() }\SpecialCharTok{\%\textgreater{}\%}
  \FunctionTok{mutate}\NormalTok{(}
    \AttributeTok{p.adj =} \FunctionTok{p.adjust}\NormalTok{(p\_value, }\AttributeTok{method =} \StringTok{"fdr"}\NormalTok{),}
    \AttributeTok{significance =} \FunctionTok{cut}\NormalTok{(p.adj, }
                       \AttributeTok{breaks =} \FunctionTok{c}\NormalTok{(}\DecValTok{0}\NormalTok{, }\FloatTok{1e{-}04}\NormalTok{, }\FloatTok{0.001}\NormalTok{, }\FloatTok{0.01}\NormalTok{, }\FloatTok{0.05}\NormalTok{, }\DecValTok{1}\NormalTok{),}
                       \AttributeTok{labels =} \FunctionTok{c}\NormalTok{(}\StringTok{"****"}\NormalTok{, }\StringTok{"***"}\NormalTok{, }\StringTok{"**"}\NormalTok{, }\StringTok{"*"}\NormalTok{, }\StringTok{"ns"}\NormalTok{))}
\NormalTok{  )}

\CommentTok{\# Display the result using kable}
\FunctionTok{kable}\NormalTok{(stat.test\_FW, }\AttributeTok{caption =} \StringTok{"Statistical Test Results for Fresh Weight (E{-}IV)"}\NormalTok{)}
\end{Highlighting}
\end{Shaded}

\begin{longtable}[]{@{}lrrl@{}}
\caption{Statistical Test Results for Fresh Weight
(E-IV)}\tabularnewline
\toprule\noalign{}
line & p\_value & p.adj & significance \\
\midrule\noalign{}
\endfirsthead
\toprule\noalign{}
line & p\_value & p.adj & significance \\
\midrule\noalign{}
\endhead
\bottomrule\noalign{}
\endlastfoot
WT & 0.0146263 & 0.0438788 & * \\
mmf1-1 & 0.4793053 & 0.4793053 & ns \\
mmf1-2 & 0.2152555 & 0.3228833 & ns \\
\end{longtable}

\begin{itemize}
\tightlist
\item
  \textbf{Purpose}: We performed ANOVA and FDR correction to assess the
  statistical significance of differences in fresh weight for MMF1
  mutants.
\end{itemize}

\subsubsection{Perform ANOVA and FDR Correction for Rosette
Area}\label{perform-anova-and-fdr-correction-for-rosette-area-1}

\begin{Shaded}
\begin{Highlighting}[]
\NormalTok{stat.test\_region\_area }\OtherTok{\textless{}{-}}\NormalTok{ phenotypes\_E2 }\SpecialCharTok{\%\textgreater{}\%}
  \FunctionTok{group\_by}\NormalTok{(line) }\SpecialCharTok{\%\textgreater{}\%}
  \FunctionTok{do}\NormalTok{(\{}
\NormalTok{    model }\OtherTok{\textless{}{-}} \FunctionTok{aov}\NormalTok{(plant\_region\_area }\SpecialCharTok{\textasciitilde{}}\NormalTok{ BX, }\AttributeTok{data =}\NormalTok{ .)}
\NormalTok{    anova\_results }\OtherTok{\textless{}{-}} \FunctionTok{summary}\NormalTok{(model)[[}\DecValTok{1}\NormalTok{]]}
\NormalTok{    p\_value }\OtherTok{\textless{}{-}}\NormalTok{ anova\_results[[}\StringTok{"Pr(\textgreater{}F)"}\NormalTok{]][}\DecValTok{1}\NormalTok{]}
    \FunctionTok{data.frame}\NormalTok{(}\AttributeTok{p\_value =}\NormalTok{ p\_value)}
\NormalTok{  \}) }\SpecialCharTok{\%\textgreater{}\%}
  \FunctionTok{ungroup}\NormalTok{() }\SpecialCharTok{\%\textgreater{}\%}
  \FunctionTok{mutate}\NormalTok{(}
    \AttributeTok{p.adj =} \FunctionTok{p.adjust}\NormalTok{(p\_value, }\AttributeTok{method =} \StringTok{"fdr"}\NormalTok{),}
    \AttributeTok{significance =} \FunctionTok{cut}\NormalTok{(}
\NormalTok{      p.adj,}
      \AttributeTok{breaks =} \FunctionTok{c}\NormalTok{(}\DecValTok{0}\NormalTok{, }\FloatTok{1e{-}04}\NormalTok{, }\FloatTok{0.001}\NormalTok{, }\FloatTok{0.01}\NormalTok{, }\FloatTok{0.05}\NormalTok{, }\DecValTok{1}\NormalTok{),}
      \AttributeTok{labels =} \FunctionTok{c}\NormalTok{(}\StringTok{"****"}\NormalTok{, }\StringTok{"***"}\NormalTok{, }\StringTok{"**"}\NormalTok{, }\StringTok{"*"}\NormalTok{, }\StringTok{"ns"}\NormalTok{)}
\NormalTok{    )}
\NormalTok{  )}

\CommentTok{\# Display the results using kable}
\FunctionTok{library}\NormalTok{(knitr)}
\FunctionTok{kable}\NormalTok{(stat.test\_region\_area, }\AttributeTok{caption =} \StringTok{"Statistical Test Results for Plant Region Area (E{-}IV)"}\NormalTok{)}
\end{Highlighting}
\end{Shaded}

\begin{longtable}[]{@{}lrrl@{}}
\caption{Statistical Test Results for Plant Region Area
(E-IV)}\tabularnewline
\toprule\noalign{}
line & p\_value & p.adj & significance \\
\midrule\noalign{}
\endfirsthead
\toprule\noalign{}
line & p\_value & p.adj & significance \\
\midrule\noalign{}
\endhead
\bottomrule\noalign{}
\endlastfoot
WT & 0.0112422 & 0.0337265 & * \\
mmf1-1 & 0.2127660 & 0.3191489 & ns \\
mmf1-2 & 0.9210433 & 0.9210433 & ns \\
\end{longtable}

\begin{itemize}
\tightlist
\item
  \textbf{Purpose}: We performed ANOVA and FDR correction to assess the
  statistical significance of differences in plant region area for MMF1
  mutants.
\end{itemize}

\section{\texorpdfstring{Experiment V: Screening of \emph{mmf1-1} and
\emph{mmf1-2} on BS09 soil
batch}{Experiment V: Screening of mmf1-1 and mmf1-2 on BS09 soil batch}}\label{experiment-v-screening-of-mmf1-1-and-mmf1-2-on-bs09-soil-batch}

\subsection{Modify and Verify Data}\label{modify-and-verify-data-2}

\begin{Shaded}
\begin{Highlighting}[]
\CommentTok{\# Change the levels of the \textquotesingle{}line\textquotesingle{} column}
\NormalTok{weight\_E3 }\OtherTok{\textless{}{-}}\NormalTok{ weight\_E3 }\SpecialCharTok{\%\textgreater{}\%}
  \FunctionTok{mutate}\NormalTok{(}\AttributeTok{genotype =} \FunctionTok{factor}\NormalTok{(genotype, }\AttributeTok{levels =} \FunctionTok{c}\NormalTok{(}\StringTok{"Col0"}\NormalTok{, }\StringTok{"mmf1{-}1"}\NormalTok{, }\StringTok{"mmf1{-}2"}\NormalTok{)))}

\CommentTok{\# Filter out rows with NA in the fresh\_weight column}
\NormalTok{weight\_E3 }\OtherTok{\textless{}{-}}\NormalTok{ weight\_E3 }\SpecialCharTok{\%\textgreater{}\%}
  \FunctionTok{filter}\NormalTok{(}\SpecialCharTok{!}\FunctionTok{is.na}\NormalTok{(fresh\_weight))}

\CommentTok{\#number of replicates}
\FunctionTok{kable}\NormalTok{(}
  \FunctionTok{table}\NormalTok{(weight\_E3}\SpecialCharTok{$}\NormalTok{soil, weight\_E3}\SpecialCharTok{$}\NormalTok{genotype), }
  \AttributeTok{caption =} \StringTok{"Number of replicates (E{-}V)"}
\NormalTok{)}
\end{Highlighting}
\end{Shaded}

\begin{longtable}[]{@{}lrrr@{}}
\caption{Number of replicates (E-V)}\tabularnewline
\toprule\noalign{}
& Col0 & mmf1-1 & mmf1-2 \\
\midrule\noalign{}
\endfirsthead
\toprule\noalign{}
& Col0 & mmf1-1 & mmf1-2 \\
\midrule\noalign{}
\endhead
\bottomrule\noalign{}
\endlastfoot
minus & 15 & 11 & 14 \\
plus & 15 & 12 & 14 \\
\end{longtable}

\begin{itemize}
\tightlist
\item
  \textbf{Purpose}: We re-leveled the `line' factor and filtered out
  rows with missing fresh weight data.
\end{itemize}

\subsection{Plot Fresh weight}\label{plot-fresh-weight-2}

\begin{Shaded}
\begin{Highlighting}[]
\CommentTok{\# Generate the boxplot}
\NormalTok{fresh\_weight\_E3 }\OtherTok{\textless{}{-}}\NormalTok{ weight\_E3 }\SpecialCharTok{\%\textgreater{}\%}
  \FunctionTok{ggplot}\NormalTok{(}\FunctionTok{aes}\NormalTok{(}\AttributeTok{x =}\NormalTok{ genotype, }\AttributeTok{y =}\NormalTok{ fresh\_weight, }\AttributeTok{fill =}\NormalTok{ soil)) }\SpecialCharTok{+} 
  \FunctionTok{geom\_boxplot}\NormalTok{(}\AttributeTok{position =} \FunctionTok{position\_dodge}\NormalTok{(}\AttributeTok{width =} \FloatTok{0.85}\NormalTok{), }\AttributeTok{alpha =} \FloatTok{0.6}\NormalTok{, }
               \AttributeTok{color =} \StringTok{"grey35"}\NormalTok{, }\AttributeTok{outlier.colour =} \ConstantTok{NA}\NormalTok{, }\AttributeTok{size =} \FloatTok{0.5}\NormalTok{) }\SpecialCharTok{+} 
\NormalTok{  ggbeeswarm}\SpecialCharTok{::}\FunctionTok{geom\_quasirandom}\NormalTok{(}\AttributeTok{dodge.width =} \FloatTok{0.85}\NormalTok{, }\AttributeTok{size =} \DecValTok{1}\NormalTok{, }\AttributeTok{shape =} \DecValTok{21}\NormalTok{, }
                               \AttributeTok{color =} \StringTok{"black"}\NormalTok{, }\AttributeTok{width =} \FloatTok{0.1}\NormalTok{, }\AttributeTok{alpha =} \FloatTok{0.4}\NormalTok{, }
                               \AttributeTok{show.legend =} \ConstantTok{FALSE}\NormalTok{) }\SpecialCharTok{+}
  \FunctionTok{stat\_summary}\NormalTok{(}\AttributeTok{fun.data =}\NormalTok{ mean\_se, }\AttributeTok{geom =} \StringTok{"errorbar"}\NormalTok{, }\AttributeTok{linewidth =} \FloatTok{0.5}\NormalTok{, }
               \AttributeTok{width =} \FloatTok{0.4}\NormalTok{, }\AttributeTok{position =} \FunctionTok{position\_dodge}\NormalTok{(}\AttributeTok{width =} \FloatTok{0.85}\NormalTok{)) }\SpecialCharTok{+} 
  \FunctionTok{stat\_compare\_means}\NormalTok{(}\AttributeTok{method =} \StringTok{"t.test"}\NormalTok{, }\AttributeTok{paired =} \ConstantTok{FALSE}\NormalTok{, }\AttributeTok{size =} \DecValTok{2}\NormalTok{, }
                     \AttributeTok{label.y =} \FloatTok{0.380}\NormalTok{, }\AttributeTok{label.x =} \DecValTok{1}\NormalTok{) }\SpecialCharTok{+}
  \FunctionTok{theme\_bw}\NormalTok{() }\SpecialCharTok{+}
  \FunctionTok{scale\_fill\_manual}\NormalTok{(}\AttributeTok{name =} \StringTok{"Conditioning"}\NormalTok{, }\AttributeTok{values =} \FunctionTok{c}\NormalTok{(}\StringTok{"seagreen"}\NormalTok{, }\StringTok{"gold2"}\NormalTok{)) }\SpecialCharTok{+} 
  \FunctionTok{coord\_cartesian}\NormalTok{(}\AttributeTok{ylim =} \FunctionTok{c}\NormalTok{(}\DecValTok{0}\NormalTok{, }\FloatTok{0.40}\NormalTok{)) }\SpecialCharTok{+}
  \FunctionTok{ggtitle}\NormalTok{(}\StringTok{"E3 {-} Fresh Weight Analysis"}\NormalTok{) }\SpecialCharTok{+}  \CommentTok{\# Add a title}
  \FunctionTok{ylab}\NormalTok{(}\StringTok{"Fresh Weight"}\NormalTok{) }\SpecialCharTok{+} 
  \FunctionTok{scale\_x\_discrete}\NormalTok{(}\AttributeTok{labels =} \FunctionTok{c}\NormalTok{(}\StringTok{"WT"}\NormalTok{, }\FunctionTok{expression}\NormalTok{(}\FunctionTok{italic}\NormalTok{(}\StringTok{"mmf1{-}1"}\NormalTok{)), }
                              \FunctionTok{expression}\NormalTok{(}\FunctionTok{italic}\NormalTok{(}\StringTok{"mmf1{-}2"}\NormalTok{)))) }\SpecialCharTok{+}  \CommentTok{\# Adjust x{-}axis labels}
  \FunctionTok{theme}\NormalTok{(}
    \AttributeTok{plot.title =} \FunctionTok{element\_text}\NormalTok{(}\AttributeTok{size =} \DecValTok{7}\NormalTok{, }\AttributeTok{hjust =} \DecValTok{0}\NormalTok{),  }\CommentTok{\# Title aligned to the left}
    \AttributeTok{axis.title.x =} \FunctionTok{element\_blank}\NormalTok{(),  }\CommentTok{\# Optional: Remove x{-}axis title}
    \AttributeTok{axis.title.y =} \FunctionTok{element\_text}\NormalTok{(}\AttributeTok{size =} \DecValTok{7}\NormalTok{),}
    \AttributeTok{axis.text.x =} \FunctionTok{element\_text}\NormalTok{(}\AttributeTok{size =} \DecValTok{7}\NormalTok{),}
    \AttributeTok{axis.text.y =} \FunctionTok{element\_text}\NormalTok{(}\AttributeTok{size =} \DecValTok{7}\NormalTok{),}
    \AttributeTok{legend.title =} \FunctionTok{element\_text}\NormalTok{(}\AttributeTok{size =} \DecValTok{7}\NormalTok{),}
    \AttributeTok{legend.text =} \FunctionTok{element\_text}\NormalTok{(}\AttributeTok{size =} \DecValTok{7}\NormalTok{),}
    \AttributeTok{strip.text =} \FunctionTok{element\_text}\NormalTok{(}\AttributeTok{size =} \DecValTok{7}\NormalTok{),  }\CommentTok{\# Facet strip text size}
    \AttributeTok{legend.position =} \StringTok{"bottom"}
\NormalTok{  )}

\CommentTok{\# Display the plot}
\FunctionTok{print}\NormalTok{(fresh\_weight\_E3)}
\CommentTok{\# Purpose: Generate and display a boxplot for fresh weight data in MMF1 mutants.}
\end{Highlighting}
\end{Shaded}

\begin{figure}[H]

\centering{

\pandocbounded{\includegraphics[keepaspectratio]{04_Mutant_feedbacks_files/figure-pdf/fig-mmf1-fresh-weight-E3-1.pdf}}

}

\caption{\label{fig-mmf1-fresh-weight-E3}Fresh weight of mmf1 mutants
(E-V).}

\end{figure}%

\subsection{Plot Plant Area for MMF1
Mutants}\label{plot-plant-area-for-mmf1-mutants-1}

\begin{Shaded}
\begin{Highlighting}[]
\CommentTok{\# Generate the boxplot}
\NormalTok{plant\_area\_E3 }\OtherTok{\textless{}{-}}\NormalTok{ phenotypes\_E3 }\SpecialCharTok{\%\textgreater{}\%}
  \FunctionTok{ggplot}\NormalTok{(}\FunctionTok{aes}\NormalTok{(}\AttributeTok{x =}\NormalTok{ genotype, }\AttributeTok{y =}\NormalTok{ plant\_region\_area, }\AttributeTok{fill =}\NormalTok{ soil)) }\SpecialCharTok{+} 
  \FunctionTok{geom\_boxplot}\NormalTok{(}\AttributeTok{position =} \FunctionTok{position\_dodge}\NormalTok{(}\AttributeTok{width =} \FloatTok{0.85}\NormalTok{), }\AttributeTok{alpha =} \FloatTok{0.6}\NormalTok{, }
               \AttributeTok{color =} \StringTok{"grey35"}\NormalTok{, }\AttributeTok{outlier.colour =} \ConstantTok{NA}\NormalTok{, }\AttributeTok{size =} \FloatTok{0.5}\NormalTok{) }\SpecialCharTok{+} 
\NormalTok{  ggbeeswarm}\SpecialCharTok{::}\FunctionTok{geom\_quasirandom}\NormalTok{(}\AttributeTok{dodge.width =} \FloatTok{0.85}\NormalTok{, }\AttributeTok{size =} \DecValTok{1}\NormalTok{, }\AttributeTok{shape =} \DecValTok{21}\NormalTok{, }
                               \AttributeTok{color =} \StringTok{"black"}\NormalTok{, }\AttributeTok{width =} \FloatTok{0.1}\NormalTok{, }\AttributeTok{alpha =} \FloatTok{0.4}\NormalTok{, }
                               \AttributeTok{show.legend =} \ConstantTok{FALSE}\NormalTok{) }\SpecialCharTok{+}
  \FunctionTok{stat\_summary}\NormalTok{(}\AttributeTok{fun.data =}\NormalTok{ mean\_se, }\AttributeTok{geom =} \StringTok{"errorbar"}\NormalTok{, }\AttributeTok{linewidth =} \FloatTok{0.5}\NormalTok{, }
               \AttributeTok{width =} \FloatTok{0.4}\NormalTok{, }\AttributeTok{position =} \FunctionTok{position\_dodge}\NormalTok{(}\AttributeTok{width =} \FloatTok{0.85}\NormalTok{)) }\SpecialCharTok{+} 
  \FunctionTok{stat\_compare\_means}\NormalTok{(}\AttributeTok{method =} \StringTok{"t.test"}\NormalTok{, }\AttributeTok{paired =} \ConstantTok{FALSE}\NormalTok{, }\AttributeTok{size =} \DecValTok{2}\NormalTok{, }
                     \AttributeTok{label.y =} \DecValTok{2800000}\NormalTok{, }\AttributeTok{label.x =} \DecValTok{1}\NormalTok{) }\SpecialCharTok{+}
  \FunctionTok{theme\_bw}\NormalTok{() }\SpecialCharTok{+}
  \FunctionTok{scale\_fill\_manual}\NormalTok{(}\AttributeTok{name =} \StringTok{"Conditioning"}\NormalTok{, }\AttributeTok{values =} \FunctionTok{c}\NormalTok{(}\StringTok{"seagreen"}\NormalTok{, }\StringTok{"gold2"}\NormalTok{)) }\SpecialCharTok{+} 
  \FunctionTok{coord\_cartesian}\NormalTok{(}\AttributeTok{ylim =} \FunctionTok{c}\NormalTok{(}\DecValTok{0}\NormalTok{, }\DecValTok{3000000}\NormalTok{)) }\SpecialCharTok{+}
  \FunctionTok{ggtitle}\NormalTok{(}\StringTok{"E3 {-} Plant Area"}\NormalTok{) }\SpecialCharTok{+}  \CommentTok{\# Add a title}
  \FunctionTok{ylab}\NormalTok{(}\StringTok{"Plant Area"}\NormalTok{) }\SpecialCharTok{+} 
  \FunctionTok{scale\_x\_discrete}\NormalTok{(}\AttributeTok{labels =} \FunctionTok{c}\NormalTok{(}\StringTok{"WT"}\NormalTok{, }\FunctionTok{expression}\NormalTok{(}\FunctionTok{italic}\NormalTok{(}\StringTok{"mmf1{-}1"}\NormalTok{)), }
                              \FunctionTok{expression}\NormalTok{(}\FunctionTok{italic}\NormalTok{(}\StringTok{"mmf1{-}2"}\NormalTok{)))) }\SpecialCharTok{+}  \CommentTok{\# Adjust x{-}axis labels}
  \FunctionTok{theme}\NormalTok{(}
    \AttributeTok{plot.title =} \FunctionTok{element\_text}\NormalTok{(}\AttributeTok{size =} \DecValTok{7}\NormalTok{, }\AttributeTok{hjust =} \DecValTok{0}\NormalTok{),  }\CommentTok{\# Title aligned to the left}
    \AttributeTok{axis.title.x =} \FunctionTok{element\_blank}\NormalTok{(),  }\CommentTok{\# Optional: Remove x{-}axis title}
    \AttributeTok{axis.title.y =} \FunctionTok{element\_text}\NormalTok{(}\AttributeTok{size =} \DecValTok{7}\NormalTok{),}
    \AttributeTok{axis.text.x =} \FunctionTok{element\_text}\NormalTok{(}\AttributeTok{size =} \DecValTok{7}\NormalTok{),}
    \AttributeTok{axis.text.y =} \FunctionTok{element\_text}\NormalTok{(}\AttributeTok{size =} \DecValTok{7}\NormalTok{),}
    \AttributeTok{legend.title =} \FunctionTok{element\_text}\NormalTok{(}\AttributeTok{size =} \DecValTok{7}\NormalTok{),}
    \AttributeTok{legend.text =} \FunctionTok{element\_text}\NormalTok{(}\AttributeTok{size =} \DecValTok{7}\NormalTok{),}
    \AttributeTok{strip.text =} \FunctionTok{element\_text}\NormalTok{(}\AttributeTok{size =} \DecValTok{7}\NormalTok{),  }\CommentTok{\# Facet strip text size}
    \AttributeTok{legend.position =} \StringTok{"bottom"}
\NormalTok{  )}

\CommentTok{\# Display the plot}
\FunctionTok{print}\NormalTok{(plant\_area\_E3)}
\CommentTok{\# Purpose: Generate and display a boxplot for plant area data in MMF1 mutants.}
\end{Highlighting}
\end{Shaded}

\begin{figure}[H]

\centering{

\pandocbounded{\includegraphics[keepaspectratio]{04_Mutant_feedbacks_files/figure-pdf/fig-mmf1-rosette-area-EV-1.pdf}}

}

\caption{\label{fig-mmf1-rosette-area-EV}Rosette area of mmf1 mutants
(E-V).}

\end{figure}%

\subsection{Statistics}\label{statistics-2}

\begin{longtable}[]{@{}lrrrrr@{}}
\caption{ANOVA Table for Fresh Weight (E-V)}\tabularnewline
\toprule\noalign{}
& Df & Sum Sq & Mean Sq & F value & Pr(\textgreater F) \\
\midrule\noalign{}
\endfirsthead
\toprule\noalign{}
& Df & Sum Sq & Mean Sq & F value & Pr(\textgreater F) \\
\midrule\noalign{}
\endhead
\bottomrule\noalign{}
\endlastfoot
genotype & 2 & 0.5130110 & 0.2565055 & 148.744250 & 0.0000000 \\
soil & 1 & 0.0081285 & 0.0081285 & 4.713604 & 0.0330885 \\
genotype:soil & 2 & 0.0058497 & 0.0029249 & 1.696092 & 0.1903569 \\
Residuals & 75 & 0.1293355 & 0.0017245 & NA & NA \\
\end{longtable}

\begin{longtable}[]{@{}lrrrrr@{}}
\caption{ANOVA Table for Plant Region Area (E-V)}\tabularnewline
\toprule\noalign{}
& Df & Sum Sq & Mean Sq & F value & Pr(\textgreater F) \\
\midrule\noalign{}
\endfirsthead
\toprule\noalign{}
& Df & Sum Sq & Mean Sq & F value & Pr(\textgreater F) \\
\midrule\noalign{}
\endhead
\bottomrule\noalign{}
\endlastfoot
genotype & 2 & 1.566506e+13 & 7.832531e+12 & 195.131941 & 0.0000000 \\
soil & 1 & 2.417528e+11 & 2.417528e+11 & 6.022790 & 0.0166883 \\
genotype:soil & 2 & 2.337858e+11 & 1.168929e+11 & 2.912154 &
0.0611680 \\
Residuals & 68 & 2.729497e+12 & 4.013967e+10 & NA & NA \\
\end{longtable}

\subsubsection{Perform ANOVA and FDR Correction for fresh
weight}\label{perform-anova-and-fdr-correction-for-fresh-weight-2}

\begin{Shaded}
\begin{Highlighting}[]
\CommentTok{\# Perform ANOVA and FDR correction for fresh weight}
\NormalTok{stat.test\_FW }\OtherTok{\textless{}{-}}\NormalTok{ weight\_E3 }\SpecialCharTok{\%\textgreater{}\%}
  \FunctionTok{group\_by}\NormalTok{(genotype) }\SpecialCharTok{\%\textgreater{}\%}
  \FunctionTok{do}\NormalTok{(\{}
\NormalTok{    model }\OtherTok{\textless{}{-}} \FunctionTok{aov}\NormalTok{(fresh\_weight }\SpecialCharTok{\textasciitilde{}}\NormalTok{ soil, }\AttributeTok{data =}\NormalTok{ .)}
\NormalTok{    anova\_results }\OtherTok{\textless{}{-}} \FunctionTok{summary}\NormalTok{(model)[[}\DecValTok{1}\NormalTok{]]}
\NormalTok{    p\_value }\OtherTok{\textless{}{-}}\NormalTok{ anova\_results[[}\StringTok{"Pr(\textgreater{}F)"}\NormalTok{]][}\DecValTok{1}\NormalTok{]}
    \FunctionTok{data.frame}\NormalTok{(}\AttributeTok{p\_value =}\NormalTok{ p\_value)}
\NormalTok{  \}) }\SpecialCharTok{\%\textgreater{}\%}
  \FunctionTok{ungroup}\NormalTok{() }\SpecialCharTok{\%\textgreater{}\%}
  \FunctionTok{mutate}\NormalTok{(}
    \AttributeTok{p.adj =} \FunctionTok{p.adjust}\NormalTok{(p\_value, }\AttributeTok{method =} \StringTok{"fdr"}\NormalTok{),}
    \AttributeTok{significance =} \FunctionTok{cut}\NormalTok{(p.adj, }
                       \AttributeTok{breaks =} \FunctionTok{c}\NormalTok{(}\DecValTok{0}\NormalTok{, }\FloatTok{1e{-}04}\NormalTok{, }\FloatTok{0.001}\NormalTok{, }\FloatTok{0.01}\NormalTok{, }\FloatTok{0.05}\NormalTok{, }\DecValTok{1}\NormalTok{),}
                       \AttributeTok{labels =} \FunctionTok{c}\NormalTok{(}\StringTok{"****"}\NormalTok{, }\StringTok{"***"}\NormalTok{, }\StringTok{"**"}\NormalTok{, }\StringTok{"*"}\NormalTok{, }\StringTok{"ns"}\NormalTok{))}
\NormalTok{  )}

\CommentTok{\# Display the result using kable}
\FunctionTok{kable}\NormalTok{(stat.test\_FW, }\AttributeTok{caption =} \StringTok{"Statistical Test Results for Fresh Weight (E{-}V)"}\NormalTok{)}
\end{Highlighting}
\end{Shaded}

\begin{longtable}[]{@{}lrrl@{}}
\caption{Statistical Test Results for Fresh Weight (E-V)}\tabularnewline
\toprule\noalign{}
genotype & p\_value & p.adj & significance \\
\midrule\noalign{}
\endfirsthead
\toprule\noalign{}
genotype & p\_value & p.adj & significance \\
\midrule\noalign{}
\endhead
\bottomrule\noalign{}
\endlastfoot
Col0 & 0.0355580 & 0.1066741 & ns \\
mmf1-1 & 0.5680571 & 0.6078448 & ns \\
mmf1-2 & 0.6078448 & 0.6078448 & ns \\
\end{longtable}

\subsubsection{Perform ANOVA and FDR Correction for Rosette
Area}\label{perform-anova-and-fdr-correction-for-rosette-area-2}

\begin{Shaded}
\begin{Highlighting}[]
\NormalTok{stat.test\_region\_area }\OtherTok{\textless{}{-}}\NormalTok{ phenotypes\_E3 }\SpecialCharTok{\%\textgreater{}\%}
  \FunctionTok{group\_by}\NormalTok{(genotype) }\SpecialCharTok{\%\textgreater{}\%}
  \FunctionTok{do}\NormalTok{(\{}
\NormalTok{    model }\OtherTok{\textless{}{-}} \FunctionTok{aov}\NormalTok{(plant\_region\_area }\SpecialCharTok{\textasciitilde{}}\NormalTok{ soil, }\AttributeTok{data =}\NormalTok{ .)}
\NormalTok{    anova\_results }\OtherTok{\textless{}{-}} \FunctionTok{summary}\NormalTok{(model)[[}\DecValTok{1}\NormalTok{]]}
\NormalTok{    p\_value }\OtherTok{\textless{}{-}}\NormalTok{ anova\_results[[}\StringTok{"Pr(\textgreater{}F)"}\NormalTok{]][}\DecValTok{1}\NormalTok{]}
    \FunctionTok{data.frame}\NormalTok{(}\AttributeTok{p\_value =}\NormalTok{ p\_value)}
\NormalTok{  \}) }\SpecialCharTok{\%\textgreater{}\%}
  \FunctionTok{ungroup}\NormalTok{() }\SpecialCharTok{\%\textgreater{}\%}
  \FunctionTok{mutate}\NormalTok{(}
    \AttributeTok{p.adj =} \FunctionTok{p.adjust}\NormalTok{(p\_value, }\AttributeTok{method =} \StringTok{"fdr"}\NormalTok{),}
    \AttributeTok{significance =} \FunctionTok{cut}\NormalTok{(}
\NormalTok{      p.adj,}
      \AttributeTok{breaks =} \FunctionTok{c}\NormalTok{(}\DecValTok{0}\NormalTok{, }\FloatTok{1e{-}04}\NormalTok{, }\FloatTok{0.001}\NormalTok{, }\FloatTok{0.01}\NormalTok{, }\FloatTok{0.05}\NormalTok{, }\DecValTok{1}\NormalTok{),}
      \AttributeTok{labels =} \FunctionTok{c}\NormalTok{(}\StringTok{"****"}\NormalTok{, }\StringTok{"***"}\NormalTok{, }\StringTok{"**"}\NormalTok{, }\StringTok{"*"}\NormalTok{, }\StringTok{"ns"}\NormalTok{)}
\NormalTok{    )}
\NormalTok{  )}

\CommentTok{\# Display the results using kable}
\FunctionTok{library}\NormalTok{(knitr)}
\FunctionTok{kable}\NormalTok{(stat.test\_region\_area, }\AttributeTok{caption =} \StringTok{"Statistical Test Results for Plant Region Area (E{-}V)"}\NormalTok{)}
\end{Highlighting}
\end{Shaded}

\begin{longtable}[]{@{}lrrl@{}}
\caption{Statistical Test Results for Plant Region Area
(E-V)}\tabularnewline
\toprule\noalign{}
genotype & p\_value & p.adj & significance \\
\midrule\noalign{}
\endfirsthead
\toprule\noalign{}
genotype & p\_value & p.adj & significance \\
\midrule\noalign{}
\endhead
\bottomrule\noalign{}
\endlastfoot
Col0 & 0.0088133 & 0.0264399 & * \\
mmf1-1 & 0.1452810 & 0.2179215 & ns \\
mmf1-2 & 0.9868226 & 0.9868226 & ns \\
\end{longtable}

\begin{Shaded}
\begin{Highlighting}[]
\CommentTok{\# Combine the three plots into one}
\NormalTok{mmf1\_mutants\_feedbacks }\OtherTok{\textless{}{-}} \FunctionTok{ggarrange}\NormalTok{(}
\NormalTok{  plant\_area\_E1, }
\NormalTok{  plant\_area\_E3, }
  \AttributeTok{nrow =} \DecValTok{1}\NormalTok{, }
  \AttributeTok{ncol =} \DecValTok{2}\NormalTok{, }
  \AttributeTok{common.legend =} \ConstantTok{TRUE}\NormalTok{,  }\CommentTok{\# Use a common legend}
  \AttributeTok{legend =} \StringTok{"bottom"}      \CommentTok{\# Place legend at the bottom}
\NormalTok{)}

\CommentTok{\# Display the combined plot}
\NormalTok{mmf1\_mutants\_feedbacks}

\CommentTok{\# Save the combined plot as a PDF}
\FunctionTok{ggsave}\NormalTok{(}
  \AttributeTok{filename =} \StringTok{"mmf1\_mutants\_rep\_feedbacks.pdf"}\NormalTok{, }
  \AttributeTok{plot =}\NormalTok{ mmf1\_mutants\_feedbacks, }
  \AttributeTok{width =} \FloatTok{17.5}\NormalTok{, }
  \AttributeTok{height =} \DecValTok{9}\NormalTok{, }
  \AttributeTok{units =} \StringTok{"cm"}\NormalTok{, }
  \AttributeTok{path =} \StringTok{"Output/figures/"}
\NormalTok{)}
\end{Highlighting}
\end{Shaded}

\begin{figure}[H]

\centering{

\pandocbounded{\includegraphics[keepaspectratio]{04_Mutant_feedbacks_files/figure-pdf/fig-mmf1-mutants-rep-1.pdf}}

}

\caption{\label{fig-mmf1-mutants-rep}Feedbacks of mmf1 mutants in repeat
experiments.}

\end{figure}%

\begin{Shaded}
\begin{Highlighting}[]
\CommentTok{\# Combine the three plots into one}
\NormalTok{mmf1\_mutants\_feedbacks\_FW }\OtherTok{\textless{}{-}} \FunctionTok{ggarrange}\NormalTok{(}
\NormalTok{  fresh\_weight\_E1, }
\NormalTok{  fresh\_weight\_E3, }
  \AttributeTok{nrow =} \DecValTok{1}\NormalTok{, }
  \AttributeTok{ncol =} \DecValTok{2}\NormalTok{, }
  \AttributeTok{common.legend =} \ConstantTok{TRUE}\NormalTok{,  }\CommentTok{\# Use a common legend}
  \AttributeTok{legend =} \StringTok{"bottom"}      \CommentTok{\# Place legend at the bottom}
\NormalTok{)}

\CommentTok{\# Display the combined plot}
\NormalTok{mmf1\_mutants\_feedbacks\_FW}

\CommentTok{\# Save the combined plot as a PDF}
\FunctionTok{ggsave}\NormalTok{(}
  \AttributeTok{filename =} \StringTok{"mmf1\_mutants\_rep\_feedbacks\_FW.pdf"}\NormalTok{, }
  \AttributeTok{plot =}\NormalTok{ mmf1\_mutants\_feedbacks\_FW, }
  \AttributeTok{width =} \FloatTok{17.5}\NormalTok{, }
  \AttributeTok{height =} \DecValTok{9}\NormalTok{, }
  \AttributeTok{units =} \StringTok{"cm"}\NormalTok{, }
  \AttributeTok{path =} \StringTok{"Output/figures/"}
\NormalTok{)}
\end{Highlighting}
\end{Shaded}

\begin{figure}[H]

\centering{

\pandocbounded{\includegraphics[keepaspectratio]{04_Mutant_feedbacks_files/figure-pdf/fig-mmf1-mutants-rep_FW-1.pdf}}

}

\caption{\label{fig-mmf1-mutants-rep_FW}Feedbacks in terms of Fresh
weight of mmf1 mutants in repeat experiments.}

\end{figure}%

\section{ETI mutants - Experiment VI}\label{eti-mutants---experiment-vi}

\subsection{Import Phenotype Table for ETI
Mutants}\label{import-phenotype-table-for-eti-mutants}

\begin{Shaded}
\begin{Highlighting}[]
\CommentTok{\# Read phenotype tables for ETI mutants}
\NormalTok{plant\_weight\_ETI }\OtherTok{\textless{}{-}} \FunctionTok{read\_excel}\NormalTok{(}\StringTok{"Input/Dataset4.xlsx"}\NormalTok{, }
                               \AttributeTok{sheet =} \StringTok{"ExpVI\_ETI\_weight"}\NormalTok{)}
\NormalTok{plant\_region\_area\_ETI }\OtherTok{\textless{}{-}} \FunctionTok{read\_excel}\NormalTok{(}\StringTok{"Input/Dataset4.xlsx"}\NormalTok{, }
                                    \AttributeTok{sheet =} \StringTok{"ExpVI\_ETI\_phenotypes"}\NormalTok{)}
\CommentTok{\# Purpose: Load the phenotype data for ETI mutants into R.}
\end{Highlighting}
\end{Shaded}

\subsection{Display Number of Replicates for ETI
Mutants}\label{display-number-of-replicates-for-eti-mutants}

\begin{Shaded}
\begin{Highlighting}[]
\CommentTok{\# Create the table}
\NormalTok{contingency\_table }\OtherTok{\textless{}{-}} \FunctionTok{table}\NormalTok{(plant\_weight\_ETI}\SpecialCharTok{$}\NormalTok{BX, plant\_weight\_ETI}\SpecialCharTok{$}\NormalTok{line)}

\CommentTok{\# Display the table with kable}
\FunctionTok{kable}\NormalTok{(contingency\_table, }\AttributeTok{caption =} \StringTok{"Number of Replicates {-} weight (E{-}VI)"}\NormalTok{)}
\end{Highlighting}
\end{Shaded}

\begin{longtable}[]{@{}lrrrr@{}}
\caption{Number of Replicates - weight (E-VI)}\tabularnewline
\toprule\noalign{}
& Col0 & eds1 & pad4 & sag101 \\
\midrule\noalign{}
\endfirsthead
\toprule\noalign{}
& Col0 & eds1 & pad4 & sag101 \\
\midrule\noalign{}
\endhead
\bottomrule\noalign{}
\endlastfoot
minus & 14 & 14 & 11 & 11 \\
plus & 14 & 13 & 14 & 12 \\
\end{longtable}

\begin{Shaded}
\begin{Highlighting}[]
\CommentTok{\# Purpose: Display the number of replicates for each line and BX condition in ETI mutants.}

\CommentTok{\# Create the table}
\NormalTok{contingency\_table }\OtherTok{\textless{}{-}} \FunctionTok{table}\NormalTok{(plant\_region\_area\_ETI}\SpecialCharTok{$}\NormalTok{BX, }
\NormalTok{                           plant\_region\_area\_ETI}\SpecialCharTok{$}\NormalTok{line)}

\CommentTok{\# Display the table with kable}
\FunctionTok{kable}\NormalTok{(contingency\_table, }\AttributeTok{caption =} \StringTok{"Number of Replicates {-} plant area (E{-}VI)"}\NormalTok{)}
\end{Highlighting}
\end{Shaded}

\begin{longtable}[]{@{}lrrrr@{}}
\caption{Number of Replicates - plant area (E-VI)}\tabularnewline
\toprule\noalign{}
& Col0 & eds1 & pad4 & sag101 \\
\midrule\noalign{}
\endfirsthead
\toprule\noalign{}
& Col0 & eds1 & pad4 & sag101 \\
\midrule\noalign{}
\endhead
\bottomrule\noalign{}
\endlastfoot
minus & 14 & 14 & 11 & 11 \\
plus & 14 & 13 & 14 & 12 \\
\end{longtable}

\begin{Shaded}
\begin{Highlighting}[]
\CommentTok{\# Export table to excel file}

\FunctionTok{write\_xlsx}\NormalTok{(plant\_weight\_ETI, }\AttributeTok{path =} \StringTok{"Output/Tables/E4\_ETI\_weight.xlsx"}\NormalTok{)}
\FunctionTok{write\_xlsx}\NormalTok{(plant\_region\_area\_ETI, }\AttributeTok{path =} \StringTok{"Output/Tables/E4\_ETI\_rosette\_area.xlsx"}\NormalTok{)}
\end{Highlighting}
\end{Shaded}

\subsection{Statistics}\label{statistics-3}

\subsubsection{Perform ANOVA for ETI
Mutants}\label{perform-anova-for-eti-mutants}

\begin{Shaded}
\begin{Highlighting}[]
\CommentTok{\# Perform ANOVA for ETI mutants}
\CommentTok{\# Fit the ANOVA model}
\NormalTok{ETI\_anova\_model\_Fw }\OtherTok{\textless{}{-}} \FunctionTok{aov}\NormalTok{(fresh\_weight }\SpecialCharTok{\textasciitilde{}}\NormalTok{ line }\SpecialCharTok{*}\NormalTok{ BX, }\AttributeTok{data =}\NormalTok{ plant\_weight\_ETI)}

\CommentTok{\# Extract the ANOVA table}
\NormalTok{anova\_table }\OtherTok{\textless{}{-}} \FunctionTok{summary}\NormalTok{(ETI\_anova\_model\_Fw)[[}\DecValTok{1}\NormalTok{]]}

\CommentTok{\# Use kable to display the table}
\FunctionTok{kable}\NormalTok{(anova\_table, }\AttributeTok{caption =} \StringTok{"ANOVA Table for Fresh Weight (E{-}VI)"}\NormalTok{)}
\end{Highlighting}
\end{Shaded}

\begin{longtable}[]{@{}lrrrrr@{}}
\caption{ANOVA Table for Fresh Weight (E-VI)}\tabularnewline
\toprule\noalign{}
& Df & Sum Sq & Mean Sq & F value & Pr(\textgreater F) \\
\midrule\noalign{}
\endfirsthead
\toprule\noalign{}
& Df & Sum Sq & Mean Sq & F value & Pr(\textgreater F) \\
\midrule\noalign{}
\endhead
\bottomrule\noalign{}
\endlastfoot
line & 3 & 6442.0308 & 2147.3436 & 0.4960015 & 0.6859221 \\
BX & 1 & 101029.3326 & 101029.3326 & 23.3361373 & 0.0000052 \\
line:BX & 3 & 342.7104 & 114.2368 & 0.0263868 & 0.9941739 \\
Residuals & 95 & 411284.2869 & 4329.3083 & NA & NA \\
\end{longtable}

\begin{Shaded}
\begin{Highlighting}[]
\CommentTok{\# Fit the ANOVA model}
\NormalTok{ETI\_anova\_model\_rosettearea }\OtherTok{\textless{}{-}} \FunctionTok{aov}\NormalTok{(plant\_region\_area }\SpecialCharTok{\textasciitilde{}}\NormalTok{ line }\SpecialCharTok{*}\NormalTok{ BX, }
                                   \AttributeTok{data =}\NormalTok{ plant\_region\_area\_ETI)}

\CommentTok{\# Extract the ANOVA table}
\NormalTok{anova\_table }\OtherTok{\textless{}{-}} \FunctionTok{summary}\NormalTok{(ETI\_anova\_model\_rosettearea)[[}\DecValTok{1}\NormalTok{]]}

\CommentTok{\# Use kable to display the table}
\FunctionTok{kable}\NormalTok{(anova\_table, }\AttributeTok{caption =} \StringTok{"ANOVA Table for Plant Region Area (E{-}VI)"}\NormalTok{)}
\end{Highlighting}
\end{Shaded}

\begin{longtable}[]{@{}lrrrrr@{}}
\caption{ANOVA Table for Plant Region Area (E-VI)}\tabularnewline
\toprule\noalign{}
& Df & Sum Sq & Mean Sq & F value & Pr(\textgreater F) \\
\midrule\noalign{}
\endfirsthead
\toprule\noalign{}
& Df & Sum Sq & Mean Sq & F value & Pr(\textgreater F) \\
\midrule\noalign{}
\endhead
\bottomrule\noalign{}
\endlastfoot
line & 3 & 1.706198e+12 & 5.687326e+11 & 1.9316359 & 0.1296917 \\
BX & 1 & 5.775147e+12 & 5.775147e+12 & 19.6146312 & 0.0000253 \\
line:BX & 3 & 7.012305e+10 & 2.337435e+10 & 0.0793883 & 0.9710410 \\
Residuals & 95 & 2.797090e+13 & 2.944306e+11 & NA & NA \\
\end{longtable}

\begin{Shaded}
\begin{Highlighting}[]
\CommentTok{\# Fit the ANOVA model}
\NormalTok{ETI\_anova\_model\_Dw }\OtherTok{\textless{}{-}} \FunctionTok{aov}\NormalTok{(dry\_weight }\SpecialCharTok{\textasciitilde{}}\NormalTok{ line }\SpecialCharTok{*}\NormalTok{ BX, }\AttributeTok{data =}\NormalTok{ plant\_weight\_ETI)}

\CommentTok{\# Extract the ANOVA table}
\NormalTok{anova\_table }\OtherTok{\textless{}{-}} \FunctionTok{summary}\NormalTok{(ETI\_anova\_model\_Dw)[[}\DecValTok{1}\NormalTok{]]}

\CommentTok{\# Use kable to display the table}
\FunctionTok{kable}\NormalTok{(anova\_table, }\AttributeTok{caption =} \StringTok{"ANOVA Table for Dry weight (E{-}VI)"}\NormalTok{)}
\end{Highlighting}
\end{Shaded}

\begin{longtable}[]{@{}lrrrrr@{}}
\caption{ANOVA Table for Dry weight (E-VI)}\tabularnewline
\toprule\noalign{}
& Df & Sum Sq & Mean Sq & F value & Pr(\textgreater F) \\
\midrule\noalign{}
\endfirsthead
\toprule\noalign{}
& Df & Sum Sq & Mean Sq & F value & Pr(\textgreater F) \\
\midrule\noalign{}
\endhead
\bottomrule\noalign{}
\endlastfoot
line & 3 & 29.266438 & 9.755479 & 0.2857264 & 0.8355970 \\
BX & 1 & 891.529645 & 891.529645 & 26.1118475 & 0.0000017 \\
line:BX & 3 & 4.622734 & 1.540911 & 0.0451315 & 0.9871912 \\
Residuals & 95 & 3243.558929 & 34.142726 & NA & NA \\
\end{longtable}

\begin{itemize}
\tightlist
\item
  \textbf{Purpose}: We performed ANOVA to analyze the effects of line
  and BX on fresh weight and plant region area in ETI mutants.
\end{itemize}

\subsubsection{Perform ANOVA and FDR Correction for ETI
Mutants}\label{perform-anova-and-fdr-correction-for-eti-mutants}

\begin{Shaded}
\begin{Highlighting}[]
\CommentTok{\# Perform ANOVA and FDR correction for fresh weight}
\NormalTok{ETI\_stat.test\_FW }\OtherTok{\textless{}{-}}\NormalTok{ plant\_weight\_ETI }\SpecialCharTok{\%\textgreater{}\%}
  \FunctionTok{group\_by}\NormalTok{(line) }\SpecialCharTok{\%\textgreater{}\%}
  \FunctionTok{do}\NormalTok{(\{}
\NormalTok{    model }\OtherTok{\textless{}{-}} \FunctionTok{aov}\NormalTok{(fresh\_weight }\SpecialCharTok{\textasciitilde{}}\NormalTok{ BX, }\AttributeTok{data =}\NormalTok{ .)}
\NormalTok{    anova\_results }\OtherTok{\textless{}{-}} \FunctionTok{summary}\NormalTok{(model)[[}\DecValTok{1}\NormalTok{]]}
\NormalTok{    p\_value }\OtherTok{\textless{}{-}}\NormalTok{ anova\_results[[}\StringTok{"Pr(\textgreater{}F)"}\NormalTok{]][}\DecValTok{1}\NormalTok{]}
    \FunctionTok{data.frame}\NormalTok{(}\AttributeTok{p\_value =}\NormalTok{ p\_value)}
\NormalTok{  \}) }\SpecialCharTok{\%\textgreater{}\%}
  \FunctionTok{ungroup}\NormalTok{() }\SpecialCharTok{\%\textgreater{}\%}
  \FunctionTok{mutate}\NormalTok{(}
    \AttributeTok{p.adj =} \FunctionTok{p.adjust}\NormalTok{(p\_value, }\AttributeTok{method =} \StringTok{"fdr"}\NormalTok{),}
    \AttributeTok{significance =} \FunctionTok{cut}\NormalTok{(}
\NormalTok{      p.adj, }
      \AttributeTok{breaks =} \FunctionTok{c}\NormalTok{(}\DecValTok{0}\NormalTok{, }\FloatTok{1e{-}04}\NormalTok{, }\FloatTok{0.001}\NormalTok{, }\FloatTok{0.01}\NormalTok{, }\FloatTok{0.05}\NormalTok{, }\DecValTok{1}\NormalTok{), }
      \AttributeTok{labels =} \FunctionTok{c}\NormalTok{(}\StringTok{"****"}\NormalTok{, }\StringTok{"***"}\NormalTok{, }\StringTok{"**"}\NormalTok{, }\StringTok{"*"}\NormalTok{, }\StringTok{"ns"}\NormalTok{)}
\NormalTok{    )}
\NormalTok{  )}

\CommentTok{\# Display fresh weight results using kable}
\FunctionTok{library}\NormalTok{(knitr)}
\FunctionTok{kable}\NormalTok{(ETI\_stat.test\_FW, }\AttributeTok{caption =} \StringTok{"Statistical Test Results for Fresh Weight (E{-}VI)"}\NormalTok{)}
\end{Highlighting}
\end{Shaded}

\begin{longtable}[]{@{}lrrl@{}}
\caption{Statistical Test Results for Fresh Weight
(E-VI)}\tabularnewline
\toprule\noalign{}
line & p\_value & p.adj & significance \\
\midrule\noalign{}
\endfirsthead
\toprule\noalign{}
line & p\_value & p.adj & significance \\
\midrule\noalign{}
\endhead
\bottomrule\noalign{}
\endlastfoot
Col0 & 0.0029714 & 0.0118855 & * \\
eds1 & 0.0310466 & 0.0456622 & * \\
pad4 & 0.0525278 & 0.0525278 & ns \\
sag101 & 0.0342466 & 0.0456622 & * \\
\end{longtable}

\begin{Shaded}
\begin{Highlighting}[]
\CommentTok{\# Perform ANOVA and FDR correction for dry weight}
\NormalTok{ETI\_stat.test\_dry\_weight }\OtherTok{\textless{}{-}}\NormalTok{ plant\_weight\_ETI }\SpecialCharTok{\%\textgreater{}\%}
  \FunctionTok{group\_by}\NormalTok{(line) }\SpecialCharTok{\%\textgreater{}\%}
  \FunctionTok{do}\NormalTok{(\{}
\NormalTok{    model }\OtherTok{\textless{}{-}} \FunctionTok{aov}\NormalTok{(dry\_weight }\SpecialCharTok{\textasciitilde{}}\NormalTok{ BX, }\AttributeTok{data =}\NormalTok{ .)}
\NormalTok{    anova\_results }\OtherTok{\textless{}{-}} \FunctionTok{summary}\NormalTok{(model)[[}\DecValTok{1}\NormalTok{]]}
\NormalTok{    p\_value }\OtherTok{\textless{}{-}}\NormalTok{ anova\_results[[}\StringTok{"Pr(\textgreater{}F)"}\NormalTok{]][}\DecValTok{1}\NormalTok{]}
    \FunctionTok{data.frame}\NormalTok{(}\AttributeTok{p\_value =}\NormalTok{ p\_value)}
\NormalTok{  \}) }\SpecialCharTok{\%\textgreater{}\%}
  \FunctionTok{ungroup}\NormalTok{() }\SpecialCharTok{\%\textgreater{}\%}
  \FunctionTok{mutate}\NormalTok{(}
    \AttributeTok{p.adj =} \FunctionTok{p.adjust}\NormalTok{(p\_value, }\AttributeTok{method =} \StringTok{"fdr"}\NormalTok{),}
    \AttributeTok{significance =} \FunctionTok{cut}\NormalTok{(}
\NormalTok{      p.adj, }
      \AttributeTok{breaks =} \FunctionTok{c}\NormalTok{(}\DecValTok{0}\NormalTok{, }\FloatTok{1e{-}04}\NormalTok{, }\FloatTok{0.001}\NormalTok{, }\FloatTok{0.01}\NormalTok{, }\FloatTok{0.05}\NormalTok{, }\DecValTok{1}\NormalTok{), }
      \AttributeTok{labels =} \FunctionTok{c}\NormalTok{(}\StringTok{"****"}\NormalTok{, }\StringTok{"***"}\NormalTok{, }\StringTok{"**"}\NormalTok{, }\StringTok{"*"}\NormalTok{, }\StringTok{"ns"}\NormalTok{)}
\NormalTok{    )}
\NormalTok{  )}

\CommentTok{\# Display dry weight results using kable}
\FunctionTok{kable}\NormalTok{(ETI\_stat.test\_dry\_weight, }\AttributeTok{caption =} \StringTok{"Statistical Test Results for Dry Weight (E{-}VI)"}\NormalTok{)}
\end{Highlighting}
\end{Shaded}

\begin{longtable}[]{@{}lrrl@{}}
\caption{Statistical Test Results for Dry Weight (E-VI)}\tabularnewline
\toprule\noalign{}
line & p\_value & p.adj & significance \\
\midrule\noalign{}
\endfirsthead
\toprule\noalign{}
line & p\_value & p.adj & significance \\
\midrule\noalign{}
\endhead
\bottomrule\noalign{}
\endlastfoot
Col0 & 0.0018612 & 0.0074447 & ** \\
eds1 & 0.0261247 & 0.0353938 & * \\
pad4 & 0.0373417 & 0.0373417 & * \\
sag101 & 0.0265453 & 0.0353938 & * \\
\end{longtable}

\begin{Shaded}
\begin{Highlighting}[]
\CommentTok{\# Perform ANOVA and FDR correction for Rosette area}
\NormalTok{ETI\_stat.test\_rosette\_area }\OtherTok{\textless{}{-}}\NormalTok{ plant\_region\_area\_ETI }\SpecialCharTok{\%\textgreater{}\%}
  \FunctionTok{group\_by}\NormalTok{(line) }\SpecialCharTok{\%\textgreater{}\%}
  \FunctionTok{do}\NormalTok{(\{}
\NormalTok{    model }\OtherTok{\textless{}{-}} \FunctionTok{aov}\NormalTok{(plant\_region\_area }\SpecialCharTok{\textasciitilde{}}\NormalTok{ BX, }\AttributeTok{data =}\NormalTok{ .)}
\NormalTok{    anova\_results }\OtherTok{\textless{}{-}} \FunctionTok{summary}\NormalTok{(model)[[}\DecValTok{1}\NormalTok{]]}
\NormalTok{    p\_value }\OtherTok{\textless{}{-}}\NormalTok{ anova\_results[[}\StringTok{"Pr(\textgreater{}F)"}\NormalTok{]][}\DecValTok{1}\NormalTok{]}
    \FunctionTok{data.frame}\NormalTok{(}\AttributeTok{p\_value =}\NormalTok{ p\_value)}
\NormalTok{  \}) }\SpecialCharTok{\%\textgreater{}\%}
  \FunctionTok{ungroup}\NormalTok{() }\SpecialCharTok{\%\textgreater{}\%}
  \FunctionTok{mutate}\NormalTok{(}
    \AttributeTok{p.adj =} \FunctionTok{p.adjust}\NormalTok{(p\_value, }\AttributeTok{method =} \StringTok{"fdr"}\NormalTok{),}
    \AttributeTok{significance =} \FunctionTok{cut}\NormalTok{(}
\NormalTok{      p.adj, }
      \AttributeTok{breaks =} \FunctionTok{c}\NormalTok{(}\DecValTok{0}\NormalTok{, }\FloatTok{1e{-}04}\NormalTok{, }\FloatTok{0.001}\NormalTok{, }\FloatTok{0.01}\NormalTok{, }\FloatTok{0.05}\NormalTok{, }\DecValTok{1}\NormalTok{), }
      \AttributeTok{labels =} \FunctionTok{c}\NormalTok{(}\StringTok{"****"}\NormalTok{, }\StringTok{"***"}\NormalTok{, }\StringTok{"**"}\NormalTok{, }\StringTok{"*"}\NormalTok{, }\StringTok{"ns"}\NormalTok{)}
\NormalTok{    )}
\NormalTok{  )}

\CommentTok{\# Display dry weight results using kable}
\FunctionTok{kable}\NormalTok{(ETI\_stat.test\_rosette\_area, }\AttributeTok{caption =} \StringTok{"Statistical Test Results for Rosette Area (E{-}VI)"}\NormalTok{)}
\end{Highlighting}
\end{Shaded}

\begin{longtable}[]{@{}lrrl@{}}
\caption{Statistical Test Results for Rosette Area
(E-VI)}\tabularnewline
\toprule\noalign{}
line & p\_value & p.adj & significance \\
\midrule\noalign{}
\endfirsthead
\toprule\noalign{}
line & p\_value & p.adj & significance \\
\midrule\noalign{}
\endhead
\bottomrule\noalign{}
\endlastfoot
Col0 & 0.0086979 & 0.0347917 & * \\
eds1 & 0.0256198 & 0.0512396 & ns \\
pad4 & 0.0803324 & 0.0803324 & ns \\
sag101 & 0.0555717 & 0.0740956 & ns \\
\end{longtable}

\begin{Shaded}
\begin{Highlighting}[]
\CommentTok{\# Purpose: Perform ANOVA and FDR correction to assess the statistical significance of differences in fresh and dry weight for ETI mutants.}
\end{Highlighting}
\end{Shaded}

\begin{itemize}
\tightlist
\item
  \textbf{Purpose}: We performed ANOVA and FDR correction to assess the
  statistical significance of differences in fresh and dry weight for
  ETI mutants.
\end{itemize}

\subsection{Plot Dry Weight for ETI
Mutants}\label{plot-dry-weight-for-eti-mutants}

\begin{figure}

\centering{

\pandocbounded{\includegraphics[keepaspectratio]{04_Mutant_feedbacks_files/figure-pdf/fig-ETI-mutants-Dw-1.pdf}}

}

\caption{\label{fig-ETI-mutants-Dw}Feedbacks of ET mutantss (E-VI).}

\end{figure}%

\begin{itemize}
\tightlist
\item
  \textbf{Purpose}: We generated a boxplot to visualize the dry weight
  data in ETI mutants.
\end{itemize}

\subsection{Plot Fresh Weight for ETI
Mutants}\label{plot-fresh-weight-for-eti-mutants}

\begin{Shaded}
\begin{Highlighting}[]
\CommentTok{\# Generate the boxplot for fresh weight}
\NormalTok{fresh\_weight\_ETI }\OtherTok{\textless{}{-}}\NormalTok{ plant\_weight\_ETI }\SpecialCharTok{\%\textgreater{}\%}
  \FunctionTok{ggplot}\NormalTok{(}\FunctionTok{aes}\NormalTok{(}\AttributeTok{x =}\NormalTok{ line, }\AttributeTok{y =}\NormalTok{ fresh\_weight, }\AttributeTok{fill =}\NormalTok{ BX)) }\SpecialCharTok{+} 
  \FunctionTok{geom\_boxplot}\NormalTok{(}\AttributeTok{position =} \FunctionTok{position\_dodge}\NormalTok{(}\AttributeTok{width =} \FloatTok{0.85}\NormalTok{), }\AttributeTok{alpha =} \FloatTok{0.6}\NormalTok{, }
               \AttributeTok{color =} \StringTok{"grey35"}\NormalTok{, }\AttributeTok{outlier.colour =} \ConstantTok{NA}\NormalTok{, }\AttributeTok{size =} \FloatTok{0.5}\NormalTok{) }\SpecialCharTok{+} 
\NormalTok{  ggbeeswarm}\SpecialCharTok{::}\FunctionTok{geom\_quasirandom}\NormalTok{(}\AttributeTok{dodge.width =} \FloatTok{0.85}\NormalTok{, }\AttributeTok{size =} \DecValTok{1}\NormalTok{, }\AttributeTok{shape =} \DecValTok{21}\NormalTok{, }
                               \AttributeTok{color =} \StringTok{"black"}\NormalTok{,}
                               \AttributeTok{width =} \FloatTok{0.1}\NormalTok{, }\AttributeTok{alpha =} \FloatTok{0.4}\NormalTok{, }
                               \AttributeTok{show.legend =} \ConstantTok{FALSE}\NormalTok{) }\SpecialCharTok{+}
  \FunctionTok{stat\_summary}\NormalTok{(}\AttributeTok{fun.data =}\NormalTok{ mean\_se, }\AttributeTok{geom =} \StringTok{"errorbar"}\NormalTok{, }\AttributeTok{linewidth =} \FloatTok{0.5}\NormalTok{, }
               \AttributeTok{width =} \FloatTok{0.4}\NormalTok{, }\AttributeTok{position =} \FunctionTok{position\_dodge}\NormalTok{(}\AttributeTok{width =} \FloatTok{0.85}\NormalTok{)) }\SpecialCharTok{+} 
  \FunctionTok{stat\_compare\_means}\NormalTok{(}\AttributeTok{method =} \StringTok{"t.test"}\NormalTok{, }\AttributeTok{paired =} \ConstantTok{FALSE}\NormalTok{, }\AttributeTok{size =} \DecValTok{2}\NormalTok{, }
                     \AttributeTok{label.y =} \DecValTok{500}\NormalTok{, }\AttributeTok{label.x =} \DecValTok{1}\NormalTok{) }\SpecialCharTok{+}
  \FunctionTok{theme\_bw}\NormalTok{() }\SpecialCharTok{+}
  \FunctionTok{scale\_fill\_manual}\NormalTok{(}\AttributeTok{name =} \StringTok{"Conditioning"}\NormalTok{, }\AttributeTok{values =} \FunctionTok{c}\NormalTok{(}\StringTok{"seagreen"}\NormalTok{, }\StringTok{"gold2"}\NormalTok{)) }\SpecialCharTok{+} 
  \FunctionTok{coord\_cartesian}\NormalTok{(}\AttributeTok{ylim =} \FunctionTok{c}\NormalTok{(}\DecValTok{0}\NormalTok{, }\DecValTok{500}\NormalTok{)) }\SpecialCharTok{+}
  \FunctionTok{ggtitle}\NormalTok{(}\StringTok{"ETI {-} Fresh Weight"}\NormalTok{) }\SpecialCharTok{+}  \CommentTok{\# Add a title}
  \FunctionTok{ylab}\NormalTok{(}\StringTok{"Fresh Weight (mg)"}\NormalTok{) }\SpecialCharTok{+} 
  \FunctionTok{scale\_x\_discrete}\NormalTok{(}\AttributeTok{labels =} \FunctionTok{c}\NormalTok{(}\StringTok{"WT"}\NormalTok{, }\FunctionTok{expression}\NormalTok{(}\FunctionTok{italic}\NormalTok{(}\StringTok{"eds1{-}12"}\NormalTok{)), }
                              \FunctionTok{expression}\NormalTok{(}\FunctionTok{italic}\NormalTok{(}\StringTok{"pad4{-}1"}\NormalTok{)), }
                              \FunctionTok{expression}\NormalTok{(}\FunctionTok{italic}\NormalTok{(}\StringTok{"sag101{-}3"}\NormalTok{)))) }\SpecialCharTok{+}  \CommentTok{\# Adjust x{-}axis labels}
  \FunctionTok{theme}\NormalTok{(}
    \AttributeTok{plot.title =} \FunctionTok{element\_text}\NormalTok{(}\AttributeTok{size =} \DecValTok{7}\NormalTok{, }\AttributeTok{hjust =} \DecValTok{0}\NormalTok{),  }\CommentTok{\# Title aligned to the left}
    \AttributeTok{axis.title.x =} \FunctionTok{element\_text}\NormalTok{(}\AttributeTok{size =} \DecValTok{7}\NormalTok{),}
    \AttributeTok{axis.title.y =} \FunctionTok{element\_text}\NormalTok{(}\AttributeTok{size =} \DecValTok{7}\NormalTok{),}
    \AttributeTok{axis.text.x =} \FunctionTok{element\_text}\NormalTok{(}\AttributeTok{size =} \DecValTok{7}\NormalTok{),}
    \AttributeTok{axis.text.y =} \FunctionTok{element\_text}\NormalTok{(}\AttributeTok{size =} \DecValTok{7}\NormalTok{),}
    \AttributeTok{legend.title =} \FunctionTok{element\_text}\NormalTok{(}\AttributeTok{size =} \DecValTok{7}\NormalTok{),}
    \AttributeTok{legend.text =} \FunctionTok{element\_text}\NormalTok{(}\AttributeTok{size =} \DecValTok{7}\NormalTok{),}
    \AttributeTok{strip.text =} \FunctionTok{element\_text}\NormalTok{(}\AttributeTok{size =} \DecValTok{7}\NormalTok{),  }\CommentTok{\# Facet strip text size}
    \AttributeTok{legend.position =} \StringTok{"bottom"}
\NormalTok{  )}

\CommentTok{\# Display the plot}
\FunctionTok{print}\NormalTok{(fresh\_weight\_ETI)}
\CommentTok{\# Purpose: Generate and display a boxplot for fresh weight data in ETI mutants.}
\end{Highlighting}
\end{Shaded}

\begin{figure}[H]

\centering{

\pandocbounded{\includegraphics[keepaspectratio]{04_Mutant_feedbacks_files/figure-pdf/fig-ETI-mutants-Fw-1.pdf}}

}

\caption{\label{fig-ETI-mutants-Fw}Feedbacks of ET mutantss (E-VI).}

\end{figure}%

\begin{itemize}
\tightlist
\item
  \textbf{Purpose}: We generated a boxplot to visualize the fresh weight
  data in ETI mutants.
\end{itemize}

\subsection{Rosette area}\label{rosette-area}

\begin{Shaded}
\begin{Highlighting}[]
\CommentTok{\# Generate the boxplot for rosette area}
\NormalTok{plant\_area\_ETI }\OtherTok{\textless{}{-}}\NormalTok{ plant\_region\_area\_ETI }\SpecialCharTok{\%\textgreater{}\%}
  \FunctionTok{ggplot}\NormalTok{(}\FunctionTok{aes}\NormalTok{(}\AttributeTok{x =}\NormalTok{ line, }\AttributeTok{y =}\NormalTok{ plant\_region\_area, }\AttributeTok{fill =}\NormalTok{ BX)) }\SpecialCharTok{+} 
  \FunctionTok{geom\_boxplot}\NormalTok{(}\AttributeTok{position =} \FunctionTok{position\_dodge}\NormalTok{(}\AttributeTok{width =} \FloatTok{0.85}\NormalTok{), }\AttributeTok{alpha =} \FloatTok{0.6}\NormalTok{, }
               \AttributeTok{color =} \StringTok{"grey35"}\NormalTok{, }\AttributeTok{outlier.colour =} \ConstantTok{NA}\NormalTok{, }\AttributeTok{size =} \FloatTok{0.5}\NormalTok{) }\SpecialCharTok{+} 
\NormalTok{  ggbeeswarm}\SpecialCharTok{::}\FunctionTok{geom\_quasirandom}\NormalTok{(}\AttributeTok{dodge.width =} \FloatTok{0.85}\NormalTok{, }\AttributeTok{size =} \DecValTok{1}\NormalTok{, }\AttributeTok{shape =} \DecValTok{21}\NormalTok{, }
                               \AttributeTok{color =} \StringTok{"black"}\NormalTok{, }\AttributeTok{width =} \FloatTok{0.1}\NormalTok{, }\AttributeTok{alpha =} \FloatTok{0.4}\NormalTok{, }
                               \AttributeTok{show.legend =} \ConstantTok{FALSE}\NormalTok{) }\SpecialCharTok{+}
  \FunctionTok{stat\_summary}\NormalTok{(}\AttributeTok{fun.data =}\NormalTok{ mean\_se, }\AttributeTok{geom =} \StringTok{"errorbar"}\NormalTok{, }\AttributeTok{linewidth =} \FloatTok{0.5}\NormalTok{, }
               \AttributeTok{width =} \FloatTok{0.4}\NormalTok{, }\AttributeTok{position =} \FunctionTok{position\_dodge}\NormalTok{(}\AttributeTok{width =} \FloatTok{0.85}\NormalTok{)) }\SpecialCharTok{+} 
  \FunctionTok{stat\_compare\_means}\NormalTok{(}\AttributeTok{method =} \StringTok{"t.test"}\NormalTok{, }\AttributeTok{paired =} \ConstantTok{FALSE}\NormalTok{, }\AttributeTok{size =} \DecValTok{2}\NormalTok{, }
                     \AttributeTok{label.y =} \DecValTok{4800000}\NormalTok{, }\AttributeTok{label.x =} \DecValTok{1}\NormalTok{) }\SpecialCharTok{+}
  \FunctionTok{theme\_bw}\NormalTok{() }\SpecialCharTok{+}
  \FunctionTok{scale\_fill\_manual}\NormalTok{(}\AttributeTok{name =} \StringTok{"Conditioning"}\NormalTok{, }\AttributeTok{values =} \FunctionTok{c}\NormalTok{(}\StringTok{"seagreen"}\NormalTok{, }\StringTok{"gold2"}\NormalTok{)) }\SpecialCharTok{+} 
  \FunctionTok{coord\_cartesian}\NormalTok{(}\AttributeTok{ylim =} \FunctionTok{c}\NormalTok{(}\DecValTok{0}\NormalTok{, }\DecValTok{5000000}\NormalTok{)) }\SpecialCharTok{+}
  \FunctionTok{ggtitle}\NormalTok{(}\StringTok{"ETI {-} Rosette Area"}\NormalTok{) }\SpecialCharTok{+}  \CommentTok{\# Add a title}
  \FunctionTok{ylab}\NormalTok{(}\StringTok{"Plant region area"}\NormalTok{) }\SpecialCharTok{+} 
  \FunctionTok{scale\_x\_discrete}\NormalTok{(}\AttributeTok{labels =} \FunctionTok{c}\NormalTok{(}\StringTok{"WT"}\NormalTok{, }\FunctionTok{expression}\NormalTok{(}\FunctionTok{italic}\NormalTok{(}\StringTok{"eds1{-}12"}\NormalTok{)), }
                              \FunctionTok{expression}\NormalTok{(}\FunctionTok{italic}\NormalTok{(}\StringTok{"pad4{-}1"}\NormalTok{)), }
                              \FunctionTok{expression}\NormalTok{(}\FunctionTok{italic}\NormalTok{(}\StringTok{"sag101{-}3"}\NormalTok{)))) }\SpecialCharTok{+}  \CommentTok{\# Adjust x{-}axis labels}
  \FunctionTok{theme}\NormalTok{(}
    \AttributeTok{plot.title =} \FunctionTok{element\_text}\NormalTok{(}\AttributeTok{size =} \DecValTok{7}\NormalTok{, }\AttributeTok{hjust =} \DecValTok{0}\NormalTok{),  }\CommentTok{\# Title aligned to the left}
    \AttributeTok{axis.title.x =} \FunctionTok{element\_text}\NormalTok{(}\AttributeTok{size =} \DecValTok{7}\NormalTok{),}
    \AttributeTok{axis.title.y =} \FunctionTok{element\_text}\NormalTok{(}\AttributeTok{size =} \DecValTok{7}\NormalTok{),}
    \AttributeTok{axis.text.x =} \FunctionTok{element\_text}\NormalTok{(}\AttributeTok{size =} \DecValTok{7}\NormalTok{),}
    \AttributeTok{axis.text.y =} \FunctionTok{element\_text}\NormalTok{(}\AttributeTok{size =} \DecValTok{7}\NormalTok{),}
    \AttributeTok{legend.title =} \FunctionTok{element\_text}\NormalTok{(}\AttributeTok{size =} \DecValTok{7}\NormalTok{),}
    \AttributeTok{legend.text =} \FunctionTok{element\_text}\NormalTok{(}\AttributeTok{size =} \DecValTok{7}\NormalTok{),}
    \AttributeTok{strip.text =} \FunctionTok{element\_text}\NormalTok{(}\AttributeTok{size =} \DecValTok{7}\NormalTok{),  }\CommentTok{\# Facet strip text size}
    \AttributeTok{legend.position =} \StringTok{"bottom"}
\NormalTok{  )}

\CommentTok{\# Display the plot}
\FunctionTok{print}\NormalTok{(plant\_area\_ETI)}
\CommentTok{\# Purpose: Generate and display a boxplot for fresh weight data in ETI mutants.}
\end{Highlighting}
\end{Shaded}

\begin{figure}[H]

\centering{

\pandocbounded{\includegraphics[keepaspectratio]{04_Mutant_feedbacks_files/figure-pdf/fig-ETI-mutants-area-1.pdf}}

}

\caption{\label{fig-ETI-mutants-area}Feedbacks of ETI mutantss (E-VI).}

\end{figure}%

\begin{Shaded}
\begin{Highlighting}[]
\FunctionTok{library}\NormalTok{(tidyverse, }\AttributeTok{quietly =} \ConstantTok{TRUE}\NormalTok{)}

\CommentTok{\# Combine the three plots into one}
\NormalTok{ETI\_mutants\_feedbacks }\OtherTok{\textless{}{-}} \FunctionTok{ggarrange}\NormalTok{(}
\NormalTok{  plant\_area\_ETI, }
\NormalTok{  dry\_weight\_ETI, }
  \AttributeTok{nrow =} \DecValTok{1}\NormalTok{, }
  \AttributeTok{ncol =} \DecValTok{2}\NormalTok{, }
  \AttributeTok{common.legend =} \ConstantTok{TRUE}\NormalTok{,  }\CommentTok{\# Use a common legend}
  \AttributeTok{legend =} \StringTok{"bottom"}      \CommentTok{\# Place legend at the bottom}
\NormalTok{)}

\CommentTok{\# Display the combined plot}
\NormalTok{ETI\_mutants\_feedbacks}

\CommentTok{\# Save the combined plot as a PDF}
\FunctionTok{ggsave}\NormalTok{(}
  \AttributeTok{filename =} \StringTok{"ETI\_mutants\_feedbacks.pdf"}\NormalTok{, }
  \AttributeTok{plot =}\NormalTok{ ETI\_mutants\_feedbacks, }
  \AttributeTok{width =} \FloatTok{16.5}\NormalTok{, }
  \AttributeTok{height =} \DecValTok{9}\NormalTok{, }
  \AttributeTok{units =} \StringTok{"cm"}\NormalTok{, }
  \AttributeTok{path =} \StringTok{"Output/figures/"}
\NormalTok{)}
\end{Highlighting}
\end{Shaded}

\begin{figure}[H]

\centering{

\pandocbounded{\includegraphics[keepaspectratio]{04_Mutant_feedbacks_files/figure-pdf/fig-ETI-mutants-1.pdf}}

}

\caption{\label{fig-ETI-mutants}Feedbacks of ETI mutants (E-VI).}

\end{figure}%

\section{Export data as Datasets}\label{export-data-as-datasets}

\begin{Shaded}
\begin{Highlighting}[]
\CommentTok{\# Load the openxlsx package}
\FunctionTok{library}\NormalTok{(openxlsx)}

\CommentTok{\# Create a new workbook}
\NormalTok{Dataset4 }\OtherTok{\textless{}{-}} \FunctionTok{createWorkbook}\NormalTok{()}

\CommentTok{\# Add sheets from mmf1 mutants and write data frames to them}
\FunctionTok{addWorksheet}\NormalTok{(Dataset4, }\StringTok{"ExpIII\_mmf1\_weight"}\NormalTok{)}
\FunctionTok{writeData}\NormalTok{(Dataset4, }\StringTok{"ExpIII\_mmf1\_weight"}\NormalTok{, weight\_E1)}
\FunctionTok{addWorksheet}\NormalTok{(Dataset4, }\StringTok{"ExpIII\_mmf1\_phenotypes"}\NormalTok{)}
\FunctionTok{writeData}\NormalTok{(Dataset4, }\StringTok{"ExpIII\_mmf1\_phenotypes"}\NormalTok{, phenotypes\_E1)}

\FunctionTok{addWorksheet}\NormalTok{(Dataset4, }\StringTok{"ExpIV\_mmf1\_weight"}\NormalTok{)}
\FunctionTok{writeData}\NormalTok{(Dataset4, }\StringTok{"ExpIV\_mmf1\_weight"}\NormalTok{, weight\_E2)}
\FunctionTok{addWorksheet}\NormalTok{(Dataset4, }\StringTok{"ExpIV\_mmf1\_phenotypes"}\NormalTok{)}
\FunctionTok{writeData}\NormalTok{(Dataset4, }\StringTok{"ExpIV\_mmf1\_phenotypes"}\NormalTok{, phenotypes\_E2)}

\FunctionTok{addWorksheet}\NormalTok{(Dataset4, }\StringTok{"ExpV\_mmf1\_weight"}\NormalTok{)}
\FunctionTok{writeData}\NormalTok{(Dataset4, }\StringTok{"ExpV\_mmf1\_weight"}\NormalTok{, weight\_E3)}
\FunctionTok{addWorksheet}\NormalTok{(Dataset4, }\StringTok{"ExpV\_mmf1\_phenotypes"}\NormalTok{)}
\FunctionTok{writeData}\NormalTok{(Dataset4, }\StringTok{"ExpV\_mmf1\_phenotypes"}\NormalTok{, phenotypes\_E3)}


\CommentTok{\# Add sheets from ETI mutants and write data frames to them}
\FunctionTok{addWorksheet}\NormalTok{(Dataset4, }\StringTok{"ExpVI\_ETI\_weight"}\NormalTok{)}
\FunctionTok{writeData}\NormalTok{(Dataset4, }\StringTok{"ExpVI\_ETI\_weight"}\NormalTok{, plant\_weight\_ETI)}
\FunctionTok{addWorksheet}\NormalTok{(Dataset4, }\StringTok{"ExpVI\_ETI\_phenotypes"}\NormalTok{)}
\FunctionTok{writeData}\NormalTok{(Dataset4, }\StringTok{"ExpVI\_ETI\_phenotypes"}\NormalTok{, plant\_region\_area\_ETI)}



\CommentTok{\# Save the workbook to an Excel file}
\FunctionTok{saveWorkbook}\NormalTok{(Dataset4, }\StringTok{"Output/Tables/Dataset4.xlsx"}\NormalTok{, }\AttributeTok{overwrite =} \ConstantTok{TRUE}\NormalTok{)}
\end{Highlighting}
\end{Shaded}

\section{Session information}\label{session-information}

\textbf{Purpose}: Print the session information, including details about
the R version, loaded packages, and the system environment. This is
useful for reproducibility, as it documents the computational
environment in which the analysis was performed.

\begin{Shaded}
\begin{Highlighting}[]
\CommentTok{\# Get session information}
\NormalTok{SI }\OtherTok{\textless{}{-}} \FunctionTok{sessionInfo}\NormalTok{()}

\CommentTok{\# Extract the attached packages and their versions}
\NormalTok{attached\_packages }\OtherTok{\textless{}{-}}\NormalTok{ SI}\SpecialCharTok{$}\NormalTok{otherPkgs}

\CommentTok{\# Print the R version}
\FunctionTok{cat}\NormalTok{(}\StringTok{"R version:"}\NormalTok{, SI}\SpecialCharTok{$}\NormalTok{R.version}\SpecialCharTok{$}\NormalTok{version.string, }\StringTok{"}\SpecialCharTok{\textbackslash{}n\textbackslash{}n}\StringTok{"}\NormalTok{)}
\end{Highlighting}
\end{Shaded}

\begin{verbatim}
R version: R version 4.4.2 (2024-10-31) 
\end{verbatim}

\begin{Shaded}
\begin{Highlighting}[]
\CommentTok{\# Print the attached packages and their versions}
\FunctionTok{cat}\NormalTok{(}\StringTok{"Attached packages:}\SpecialCharTok{\textbackslash{}n}\StringTok{"}\NormalTok{)}
\end{Highlighting}
\end{Shaded}

\begin{verbatim}
Attached packages:
\end{verbatim}

\begin{Shaded}
\begin{Highlighting}[]
\FunctionTok{invisible}\NormalTok{(}\FunctionTok{lapply}\NormalTok{(attached\_packages, }\ControlFlowTok{function}\NormalTok{(pkg) \{}
  \FunctionTok{cat}\NormalTok{(pkg}\SpecialCharTok{$}\NormalTok{Package, }\StringTok{"{-}"}\NormalTok{, pkg}\SpecialCharTok{$}\NormalTok{Version, }\StringTok{"}\SpecialCharTok{\textbackslash{}n}\StringTok{"}\NormalTok{)}
\NormalTok{\}))}
\end{Highlighting}
\end{Shaded}

\begin{verbatim}
readxl - 1.4.3 
openxlsx - 4.2.7.1 
writexl - 1.5.1 
knitr - 1.49 
cowplot - 1.1.3 
pander - 0.6.5 
ggpubr - 0.6.0 
ggtext - 0.1.2 
lubridate - 1.9.4 
forcats - 1.0.0 
stringr - 1.5.1 
dplyr - 1.1.4 
purrr - 1.0.2 
readr - 2.1.5 
tidyr - 1.3.1 
tibble - 3.2.1 
tidyverse - 2.0.0 
ggbeeswarm - 0.7.2 
ggplot2 - 3.5.1 
\end{verbatim}




\end{document}
